\documentclass[12pt,a4paper,twoside]{article}

\usepackage[utf8]{inputenc}
\usepackage[T1]{fontenc}
\usepackage[ngerman]{babel}
\usepackage{lmodern}
\usepackage{microtype}
\usepackage{amsmath,amssymb,amsthm,mathtools}
\usepackage{geometry}
\usepackage{graphicx}
\usepackage{xcolor}
\usepackage{booktabs}
\usepackage{array}
\usepackage{longtable}
\usepackage{enumitem}
\usepackage{mdframed}
\usepackage{tikz}
\usetikzlibrary{graphs,graphs.standard,positioning,arrows.meta,
                decorations.pathmorphing,calc,shapes.geometric,
                matrix,fit,backgrounds}
\usepackage{pgfplots}
\pgfplotsset{compat=1.18}
\usepackage{algorithm}
\usepackage{algpseudocode}
\usepackage{listings}
\usepackage[hidelinks]{hyperref}
\usepackage{cleveref}

\setlist[itemize,1]{label=--}

% ---------- Farben ----------
\definecolor{mainblue}{RGB}{25, 70, 170}
\definecolor{accentred}{RGB}{180, 50, 30}
\definecolor{lightgray}{RGB}{245, 245, 248}
\definecolor{darkgray}{RGB}{60, 60, 70}
\definecolor{darkgreen}{RGB}{30, 100, 50}

\hypersetup{
	colorlinks=true,
	linkcolor=mainblue,
	citecolor=accentred,
	urlcolor=mainblue,
	pdftitle={Logistische Regression: Theorie, Algorithmen und Zukunftsperspektiven},
	pdfauthor={Stephan Epp},
}

%------------------------------------------------------------
% Cleveref – vollständig deutsche Bezeichnungen
%------------------------------------------------------------
% Benutzerdefinierte Theorem-Umgebungen
\crefname{theorem}{Satz}{Sätze}
\Crefname{theorem}{Satz}{Sätze}
\crefname{lemma}{Lemma}{Lemmata}
\Crefname{lemma}{Lemma}{Lemmata}
\crefname{korollar}{Korollar}{Korollare}
\Crefname{korollar}{Korollar}{Korollare}
\crefname{proposition}{Proposition}{Propositionen}
\Crefname{proposition}{Proposition}{Propositionen}
\crefname{definition}{Definition}{Definitionen}
\Crefname{definition}{Definition}{Definitionen}
\crefname{beispiel}{Beispiel}{Beispiele}
\Crefname{beispiel}{Beispiel}{Beispiele}
\crefname{bemerkung}{Bemerkung}{Bemerkungen}
\Crefname{bemerkung}{Bemerkung}{Bemerkungen}
\crefname{beobachtung}{Beobachtung}{Beobachtungen}
\Crefname{beobachtung}{Beobachtung}{Beobachtungen}
% Standard-Umgebungen
\crefname{figure}{Abb.}{Abb.}
\Crefname{figure}{Abbildung}{Abbildungen}
\crefname{table}{Tab.}{Tab.}
\Crefname{table}{Tabelle}{Tabellen}
\crefname{section}{Abschn.}{Abschn.}
\Crefname{section}{Abschnitt}{Abschnitte}
\crefname{equation}{Gl.}{Gl.}
\Crefname{equation}{Gleichung}{Gleichungen}
\crefname{algorithm}{Alg.}{Alg.}
\Crefname{algorithm}{Algorithmus}{Algorithmen}
\usepackage{caption}
\usepackage{subcaption}
\usepackage{setspace}
\usepackage{fancyhdr}
\usepackage{titlesec}
\usepackage{abstract}
%\usepackage{natbib}

%------------------------------------------------------------
% Seitenränder
%------------------------------------------------------------
\geometry{margin=2.5cm}

%------------------------------------------------------------
% Theorem-Umgebungen
%------------------------------------------------------------
\theoremstyle{definition}
\newtheorem{definition}{Definition}[section]
\newtheorem{beispiel}{Beispiel}[section]

\theoremstyle{plain}
\newtheorem{theorem}{Satz}[section]
\newtheorem{lemma}[theorem]{Lemma}
\newtheorem{korollar}[theorem]{Korollar}
\newtheorem{proposition}[theorem]{Proposition}

\theoremstyle{remark}
\newtheorem{bemerkung}{Bemerkung}[section]
\newtheorem{beobachtung}{Beobachtung}[section]

%------------------------------------------------------------
% Hervorgehobene Box
%------------------------------------------------------------
\newmdenv[
  backgroundcolor=gray!10,
  linecolor=black!40,
  linewidth=0.8pt,
  innerleftmargin=10pt,
  innerrightmargin=10pt,
  innertopmargin=8pt,
  innerbottommargin=8pt,
  skipabove=10pt,
  skipbelow=10pt
]{infobox}

%------------------------------------------------------------
% Listings
%------------------------------------------------------------
\lstset{
  basicstyle=\ttfamily\small,
  keywordstyle=\bfseries\color{blue!70!black},
  commentstyle=\color{green!50!black},
  stringstyle=\color{red!70!black},
  numbers=left,
  numberstyle=\tiny\color{gray},
  frame=single,
  breaklines=true,
  captionpos=b
}


\title{\textbf{Die Notwendigkeit der \\Modellierung in Graphen}}
\author{Stephan Epp}
\date{19. März 2026}

%============================================================
\begin{document}
%============================================================

\maketitle

\tableofcontents
\newpage

%============================================================
\section{Einleitung}
\label{sec:einleitung}
%============================================================
Graphen gehören zu den mächtigsten und vielseitigsten Abstraktionen,
die Mathematik und Informatik hervorgebracht haben.  Diese Arbeit
stellt die These auf – und beweist sie –, dass die Modellierung
mittels Graphen nicht nur eine nützliche Methode ist, sondern die
\emph{mathematisch optimale} Modellierungsform für alle Strukturen
ist, die als Graph darstellbar sind.  Kernelement des Beweises ist
die fundamentale Atomizität des Graphen: Jede Graphstruktur besteht
aus der kleinsten denkbaren diskreten Einheit, nämlich
\emph{zwei Knoten und einer Kante}.  Diese Binärteilung erzeugt
eine natürliche logarithmische Tiefenstruktur: Divide-and-Conquer
halbiert bei jedem Schritt den Problemraum und führt so zu
$\mathcal{O}(\log n)$-Verhalten in Tiefe, Suche und Optimierung.
Wir entwickeln diesen Universalitätsbeweis formal und illustrieren
ihn anhand ausführlicher Beispiele aus sechs wissenschaftlichen
Disziplinen: Mathematik, Genetik, Neurowissenschaft, Regelungs\-technik,
Finanzwesen und Informatik.  Als krönenden Abschluss wird der
\emph{Optimale Subgraph Algorithmus} vollständig hergeleitet,
der diese logarithmische Struktur ausnutzt, um in beliebig großen
Graphen optimale Teilstrukturen effizient zu identifizieren.

\textbf{Schlüsselwörter}: Graphentheorie, logarithmische Modellierung, optimaler Subgraph,
Divide-and-Conquer, universelle Strukturabstraktion,
Komplexitätstheorie, interdisziplinäre Modellierung.

\subsection{Motivation und Problemstellung}

Seit den Anfängen der Mathematik steht die Menschheit vor der
Herausforderung, komplexe reale Systeme in formalen Strukturen
zu erfassen.  Vom babylonischen Straßennetz über das Neuronengefüge
des menschlichen Gehirns bis hin zu Proteinfaltungsnetzwerken und
globalen Finanzmärkten – stets treten dieselben strukturellen
Grundmuster auf: diskrete \emph{Einheiten} (Knoten) und
\emph{Beziehungen} zwischen ihnen (Kanten).

Leonhard Euler legte 1736 mit der Lösung des Königsberger
Brückenproblems den Grundstein der Graphentheorie.  Was Euler
als spielerische Übung begann, hat sich zur vielleicht
universellsten Modellierungssprache der Wissenschaften entwickelt.
Dennoch fehlt bislang eine kohärente, disziplinübergreifende
Begründung, weshalb die Graphenmodellierung nicht nur
\emph{brauchbar}, sondern \emph{optimal} ist.

\subsection{Zentrale Aussage}

Jede Struktur, die als Graph modelliert werden
kann, wird durch die Graphenmodellierung \emph{optimal} repräsentiert,
weil die Elementareinheit des Graphen – \emph{zwei Knoten und eine
Kante} – eine inhärente logarithmische Tiefenstruktur erzeugt.
Diese logarithmische Struktur garantiert die informationstheoretisch
minimale Beschreibungskomplexität und die algorithmisch optimale
Laufzeit für alle wesentlichen Operationen: Suche, Optimierung,
Partitionierung und Kommunikation.

%============================================================
\section{Grundlagen}
\label{sec:grundlagen}
%============================================================

\begin{definition}[Graph]
Ein \emph{Graph} $G = (V, E)$ besteht aus einer endlichen oder
abzählbar unendlichen Menge $V$ von \emph{Knoten} (Vertices) und
einer Menge $E \subseteq \binom{V}{2}$ von \emph{Kanten} (Edges),
wobei jede Kante einem ungeordneten Paar $\{u,v\}$ von Knoten
entspricht.  Im Fall gerichteter Graphen (\emph{Digraphen}) ist
$E \subseteq V \times V$ eine Menge geordneter Paare.
\end{definition}

\begin{definition}[Gewichteter Graph]
Ein \emph{gewichteter Graph} $(G, w)$ ist ein Graph $G=(V,E)$
zusammen mit einer Gewichtsfunktion $w: E \to \mathbb{R}$.
\end{definition}

\begin{definition}[Grad, Weg, Zusammenhang]
Der \emph{Grad} $\deg(v)$ eines Knotens $v \in V$ ist die Anzahl
seiner Nachbarn.  Ein \emph{Weg} der Länge $k$ ist eine Folge
$v_0, e_1, v_1, \ldots, e_k, v_k$ von paarweise verschiedenen
Knoten.  $G$ heißt \emph{zusammenhängend}, wenn zwischen je zwei
Knoten ein Weg existiert.
\end{definition}

\begin{definition}[Baum und Spannbaum]
Ein zusammenhängender, kreisfreier Graph heißt \emph{Baum}.
Ein \emph{Spannbaum} von $G$ ist ein Teilgraph von $G$, der alle
Knoten enthält und selbst ein Baum ist.
\end{definition}

\subsection{Die Elementareinheit: Zwei Knoten, eine Kante}

Die kleinste nicht-triviale Graphstruktur ist der \emph{elementare
Graph}:
\[
  G_{\text{atom}} = \bigl(\{u, v\},\; \{(u,v)\}\bigr).
\]
Er besteht aus exakt zwei Knoten und einer Kante.  Diese Einheit
kodiert eine binäre Beziehung – das kleinste mögliche Stück
struktureller Information.  Jeder beliebig komplexe Graph entsteht
durch Komposition solcher Atome.

\begin{figure}[h]
  \centering
  \begin{tikzpicture}[
    every node/.style={circle,draw,minimum size=0.7cm,thick},
    >=Stealth
  ]
    \node (u) at (0,0) {$u$};
    \node (v) at (3,0) {$v$};
    \draw[->,thick] (u) -- node[midway,above,draw=none] {$e$} (v);
  \end{tikzpicture}
  \caption{Die Elementareinheit des Graphen: Zwei Knoten $u, v$
           und eine Kante $e=(u,v)$.  Alles beginnt hier.}
  \label{fig:atom}
\end{figure}

\subsection{Sukzessive Komposition und Baumstruktur}

Wenn wir die Elementareinheit rekursiv verdoppeln – d.h. jeden
Knoten durch ein weiteres Atom ersetzen – entsteht ein binärer
Baum.  Bei $k$ Verdoppelungsschritten ergibt sich ein Baum mit
$2^k$ Blättern und Höhe~$k$.  Umgekehrt:  Um $n$ Elemente in
einem binären Baum zu organisieren, genügen
\[
  k = \lceil \log_2 n \rceil
\]
Ebenen.  Darin liegt der Kern der logarithmischen Modellierung.

\begin{figure}[h]
  \centering
  \begin{tikzpicture}[
    level distance=1.5cm,
    level 1/.style={sibling distance=4cm},
    level 2/.style={sibling distance=2cm},
    every node/.style={circle,draw,minimum size=0.6cm,thick,font=\small}
  ]
  \node {$r$}
    child { node {$a$}
      child { node {$c$} }
      child { node {$d$} }
    }
    child { node {$b$}
      child { node {$e$} }
      child { node {$f$} }
    };
  \end{tikzpicture}
  \caption{Binärer Baum als sukzessive Komposition der
           Elementareinheit.  Tiefe $k=2$ kodiert $2^2=4$ Blätter.
           Jede Ebene halbiert den Problemraum.}
  \label{fig:binaer}
\end{figure}

%============================================================
\section{Graphen als logarithmisch optimale Modellierung}
\label{sec:log_beweis}
%============================================================

\subsection{Informationstheoretische Grundlage}

\begin{theorem}[Logarithmische Beschreibungstiefe]
\label{thm:log}
Sei $\mathcal{S}$ eine Menge von $n$ diskreten Elementen mit
einer beliebigen binären Beziehungsstruktur.  Dann ist die
minimale Anzahl von Bits, die benötigt wird, um ein beliebiges
Element eindeutig zu identifizieren, gleich
\[
  H_{\min} = \lceil \log_2 n \rceil.
\]
Diese untere Schranke wird genau dann erreicht, wenn die Struktur
als vollständiger Binärbaum – und damit als spezieller Graph –
organisiert ist.
\end{theorem}

\begin{proof}
Nach Shannons Quellencodierungstheorem~\cite{Shannon1948} benötigt
jede verlustfreie Kodierung von $n$ gleichwahrscheinlichen Symbolen
mindestens $\log_2 n$ Bits pro Symbol.  Eine Binärbaumstruktur
der Tiefe $\lceil \log_2 n \rceil$ stellt genau diese Schranke
bereit: Jeder Pfad von der Wurzel zu einem Blatt kodiert eine
Bitfolge der Länge $\lceil \log_2 n \rceil$.  Kein anderes
Kodierungsschema kann diese Schranke unterschreiten.
Da jeder Binärbaum ein spezieller Graph ist, ist die
Graphstruktur informationstheoretisch optimal.
\end{proof}

\subsection{Algorithmische Optimalität: Divide-and-Conquer}

Die Grundstrategie des Divide-and-Conquer ist strukturell identisch
mit dem Durchlaufen eines Binärbaums:

\begin{enumerate}[label=\textbf{(\arabic*)}]
  \item \textbf{Divide:} Teile das Problem $P$ der Größe $n$ in
        zwei Teilprobleme $P_L$ und $P_R$ der Größe $n/2$.
  \item \textbf{Conquer:} Löse $P_L$ und $P_R$ rekursiv.
  \item \textbf{Combine:} Füge die Teillösungen zur Gesamtlösung
        zusammen.
\end{enumerate}

\noindent
Die zugehörige Rekurrenz lautet:
\[
  T(n) = 2\,T\!\left(\tfrac{n}{2}\right) + \mathcal{O}(n).
\]
Nach dem Master-Theorem~\cite{Cormen2009} ergibt sich
$T(n) = \mathcal{O}(n \log n)$.  Bei reiner Suche (ohne
Combine-Overhead) folgt $T(n)=\mathcal{O}(\log n)$.

\begin{theorem}[Divide-and-Conquer-Optimalität]
\label{thm:dc}
Jeder Suchalgorithmus über $n$ unstrukturierten Elementen
benötigt im schlechtesten Fall $\Omega(\log n)$ Vergleiche.
Eine binäre Suchstruktur (Binärsuchbaum, AVL-Baum, B-Baum)
erreicht diese Schranke.  Da alle diese Strukturen spezielle
Graphen sind, ist die Graphmodellierung algorithmisch optimal.
\end{theorem}

\begin{proof}
Jeder Vergleich teilt die Menge verbleibender Kandidaten maximal
in zwei gleich große Hälften.  Nach $k$ Vergleichen verbleiben
höchstens $n/2^k$ Kandidaten.  Für $n/2^k \leq 1$ folgt
$k \geq \log_2 n$.  Dies ist eine informationstheoretische
untere Schranke, die unabhängig vom Algorithmus gilt.
Binäre Suchbäume erreichen genau $\mathcal{O}(\log n)$ pro
Operation.
\end{proof}

\subsection{Strukturelle Überlegenheit über lineare und
            polynomielle Modellierungen}

\begin{theorem}[Überlegenheit der logarithmischen Modellierung]
\label{thm:superior}
Für jede Strukturgröße $n \geq 3$ gilt:
\[
  \log_2 n \;<\; \sqrt{n} \;<\; n \;<\; n^2 \;<\; 2^n.
\]
Eine logarithmische Suchtiefe (realisiert durch Graphstrukturen)
ist daher asymptotisch und praktisch jeder linearen oder
polynomiellen Alternative überlegen.
\end{theorem}

\noindent
\Cref{fig:komplexitaet} veranschaulicht diesen Unterschied.

\begin{figure}[h]
  \centering
  \begin{tikzpicture}
    \begin{axis}[
      width=10cm, height=7cm,
      xlabel={Problemgröße $n$},
      ylabel={Operationen},
      xmin=1, xmax=64,
      ymin=0, ymax=70,
      legend pos=north west,
      grid=major,
      grid style={line width=0.2pt, draw=gray!30},
      thick
    ]
      \addplot[blue,  domain=1:64, samples=200] {ln(x)/ln(2)};
      \addlegendentry{$\log_2 n$}
      \addplot[green!60!black, domain=1:64, samples=200] {sqrt(x)};
      \addlegendentry{$\sqrt{n}$}
      \addplot[orange, domain=1:64, samples=200] {x};
      \addlegendentry{$n$}
      \addplot[red,    domain=1:8,  samples=100] {x^2};
      \addlegendentry{$n^2$}
    \end{axis}
  \end{tikzpicture}
  \caption{Vergleich der Wachstumsraten.  Die logarithmische Kurve
           (blau) bleibt bei wachsendem $n$ weit unter allen anderen
           Komplexitätsklassen.}
  \label{fig:komplexitaet}
\end{figure}

\subsection{Universalitätssatz}

\begin{theorem}[Universalität der Graphmodellierung]
\label{thm:universal}
Jede endliche Relationenstruktur $\mathcal{R} = (D, R_1, \ldots, R_k)$
über einer Domäne $D$ mit binären Relationen $R_i \subseteq D \times D$
ist isomorph zu einem gewichteten Digraphen $G=(V,E,w)$ mit
$|V|=|D|$ und $|E|=\sum_i |R_i|$.
\end{theorem}

\begin{proof}
Setze $V := D$, $E := \bigcup_i R_i$ und $w(u,v) := i$ für
$(u,v) \in R_i$.  Die Abbildung $\phi: d \mapsto d$ ist ein
Isomorphismus zwischen $\mathcal{R}$ und $(G,w)$.
\end{proof}

\noindent
Dieser Satz zeigt: Jede binäre Relation – und damit jede
strukturelle Abhängigkeit in Wissenschaft und Technik – ist
exakt durch einen Graphen modellierbar.  Die Frage ist nicht
\emph{ob}, sondern nur \emph{welcher} Graph.

\subsection{Das logarithmische Prinzip: Zusammenfassung}

\textbf{Das logarithmische Prinzip der Graphenmodellierung:}
\begin{enumerate}
  \item Die Atomeinheit des Graphen ($u \to v$) ist eine binäre
        Relation – das informationstheoretische Minimum.
  \item Durch Komposition dieser Atome entstehen Baumstrukturen,
        die automatisch logarithmische Tiefe haben.
  \item Divide-and-Conquer auf Graphen halbiert den Problemraum
        mit jeder Kante auf dem Suchpfad.
  \item Jede andere Modellierung (Liste, Matrix, Tensor) kann
        entweder in eine Graphstruktur konvertiert werden
        (und verliert damit nichts) oder ist schlechter.
  \item Die logarithmische Schranke $\mathcal{O}(\log n)$
        ist informationstheoretisch optimal und wird durch
        balancierende Graphstrukturen stets erreicht.
\end{enumerate}

%============================================================
\section{Anwendungen in der Mathematik}
\label{sec:math}
%============================================================

\subsection{Zahlentheorie und Primzahlgraphen}

Zahlentheoretische Strukturen zeigen auf natürliche Weise
Graphcharakter. Sei $V = \{2, 3, \ldots, n\}$ die Menge der
natürlichen Zahlen.  Wir definieren den \emph{Teilbarkeitsgraphen}:
\[
  E = \{(a,b) \mid a \text{ teilt } b,\; a \neq b\}.
\]

\begin{beispiel}[Teilbarkeitsgraph]
Für $n=12$ entsteht ein DAG, dessen Wurzeln die Primzahlen sind.
Die Primzahlzerlegung $12 = 2^2 \cdot 3$ entspricht einem Pfad
$2 \to 4 \to 12$ und $3 \to 12$ im Graphen.  Die
Fundamentalsatz-Eigenschaft der eindeutigen Primfaktorzerlegung
ist äquivalent zur Existenz eines eindeutigen minimalen
aufspannenden Waldes.
\end{beispiel}

\begin{figure}[h]
  \centering
  \begin{tikzpicture}[
    every node/.style={circle,draw,minimum size=0.65cm,thick,font=\footnotesize},
    >=Stealth,
    node distance=1.4cm
  ]
    \node (p2)  {2};
    \node (p3)  [right=2cm of p2]  {3};
    \node (p5)  [right=2cm of p3]  {5};
    \node (p7)  [right=2cm of p5]  {7};
    \node (n4)  [below=1.5cm of p2] {4};
    \node (n6)  [below=1.5cm of p3] {6};
    \node (n9)  [below=1.5cm of p3, xshift=1.5cm] {9};
    \node (n12) [below=1.5cm of n6] {12};
    \draw[->] (p2) -- (n4);
    \draw[->] (p2) -- (n6);
    \draw[->] (p3) -- (n6);
    \draw[->] (p3) -- (n9);
    \draw[->] (n4) -- (n12);
    \draw[->] (n6) -- (n12);
  \end{tikzpicture}
  \caption{Teilbarkeitsgraph der Zahlen $\{2,3,4,5,6,7,9,12\}$.
           Primzahlen sind Quellen (kein eingehender Pfeil).
           Die Graphstruktur kodiert die Primfaktorzerlegung
           vollständig.}
  \label{fig:divisibility}
\end{figure}

\subsection{Algebraische Strukturen als Cayley-Graphen}

Jede endliche Gruppe $(G, \cdot)$ mit Erzeugendenmenge $S$ besitzt
einen \emph{Cayley-Graphen} $\text{Cay}(G,S)$, der die algebraische
Struktur vollständig kodiert.

\begin{theorem}[Cayley, 1878]
Der Cayley-Graph $\text{Cay}(G,S)$ ist regulär vom Grad $|S|$.  Er
ist zusammenhängend genau dann, wenn $S$ die Gruppe $G$ erzeugt.
Der Durchmesser des Cayley-Graphen misst die maximale Wortlänge
in den Erzeugenden – ein direktes Maß für
Kommunikationseffizienz in der Gruppe.
\end{theorem}

\begin{beispiel}[Cayley-Graph von $\mathbb{Z}_8$]
Für $G = \mathbb{Z}_8$ mit $S=\{1\}$ entsteht ein 8-Kreis.
Mit $S=\{1,4\}$ entsteht ein Hyperwürfel, dessen Durchmesser
nur $\lceil\log_2 8\rceil = 3$ beträgt – logarithmisch in $|G|$.
Dies ist kein Zufall: Hyperwürfelgraphen haben stets logarithmischen
Durchmesser und sind damit für parallele Kommunikation optimal.
\end{beispiel}

\subsection{Lineare Algebra: Laplace-Operator und Spektrum}

Die \emph{Laplace-Matrix} eines Graphen $G=(V,E)$ lautet:
\[
  L = D - A,
\]
wobei $D = \operatorname{diag}(\deg(v_1), \ldots, \deg(v_n))$
die Gradmatrix und $A$ die Adjazenzmatrix ist.  Das
Spektrum $0 = \lambda_1 \leq \lambda_2 \leq \cdots \leq \lambda_n$
von $L$ trägt tiefe strukturelle Information:

\begin{itemize}
  \item $\lambda_1 = 0$: jeder Graph hat mindestens einen Nulleigenwert.
  \item Die Multiplizität von $\lambda_1 = 0$ gibt die Anzahl
        zusammenhängender Komponenten an.
  \item $\lambda_2 > 0$ (Algebraische Konnektivität / Fiedler-Wert)
        ist ein Maß für die Robustheit des Graphen.
  \item $\lambda_2$ bestimmt die Mischzeit einer Irrfahrt:
        $\tau_{\text{mix}} = \mathcal{O}(\lambda_2^{-1} \log n)$ –
        erneut logarithmisch.
\end{itemize}

\begin{beispiel}[Spektrum eines Pfadgraphen]
Der Pfadgraph $P_n$ auf $n$ Knoten hat Eigenwerte
$\lambda_k = 2 - 2\cos\!\left(\tfrac{(k-1)\pi}{n}\right)$.
Für große $n$ ergibt sich $\lambda_2 \approx \pi^2/n^2$,
was einer Mischzeit von $\mathcal{O}(n^2 \log n)$ entspricht –
ein Hinweis, dass Pfadgraphen schlecht strukturiert (lang und
linear) sind.  Expandergraphen mit $\lambda_2 = \Omega(1/\log n)$
erreichen hingegen $\mathcal{O}(\log^2 n)$ Mischzeit.
\end{beispiel}

\subsection{Topologie: Simplizialkomplexe und Graphhomologie}

Graphen bilden die $1$-Skelette simplizialer Komplexe.  Die
\emph{Betti-Zahlen} $\beta_0, \beta_1$ eines Graphen zählen
Zusammenhangskomponenten ($\beta_0$) und unabhängige Zyklen
($\beta_1 = |E| - |V| + \beta_0$, der Zyklomatische Rang).
Dies verbindet Graphenstruktur unmittelbar mit topologischen
Invarianten und erlaubt die Analyse globaler Formen komplexer
Netzwerke (Löcher, Tunnel) durch rein lokale Kanteninformation.

\subsection{Kombinatorik: Matchings und Heiratstheoreme}

\begin{theorem}[Heiratssatz, Hall 1935]
Ein bipartiter Graph $G = (A \cup B, E)$ besitzt genau dann
ein perfektes Matching, wenn für jede Teilmenge $S \subseteq A$
gilt: $|N(S)| \geq |S|$, wobei $N(S)$ die Nachbarschaft
von $S$ in $B$ bezeichnet.
\end{theorem}

Dieser Satz ist die Grundlage für Zuordnungsprobleme aller Art:
Stellenvermittlung, Aufgabenzuweisung, Ressourcenplanung.  Der
optimale Matching-Algorithmus (Ungarische Methode) arbeitet in
$\mathcal{O}(n^3)$, das Hopcroft-Karp-Verfahren auf bipartiten
Graphen sogar in $\mathcal{O}(|E| \sqrt{|V|})$.  Die logarithmische
Schichtung (BFS in $\mathcal{O}(\log n)$ Phasen) ist dabei der
Schlüssel zum Effizienzgewinn.

%============================================================
\section{Anwendungen in der Genetik}
\label{sec:genetik}
%============================================================

\subsection{Das Proteom als Interaktionsgraph}

Das menschliche Proteom umfasst etwa 20.000 Proteine, die
untereinander ein dichtes Netz von Wechselwirkungen bilden
(Protein-Protein-Interaktionsnetzwerk, PPI).  Formal:
\[
  G_{\text{PPI}} = (V_P, E_P), \quad
  V_P = \text{Proteine}, \quad
  E_P = \text{physikalische Interaktionen}.
\]

\begin{beobachtung}
PPI-Netzwerke sind skalenfreie Graphen~\cite{Barabasi1999} mit
Gradverteilung $P(\deg = k) \propto k^{-\gamma}$, $\gamma \approx 2.5$.
Die Hubs (hochgradige Knoten) sind essentielle Proteine – ihr
Knockout ist häufig letal.  Diese Erkenntnis ist nur durch
Graphmodellierung möglich, nicht durch isolierte Proteinanalyse.
\end{beobachtung}

\begin{beispiel}[Hub-Protein p53]
Das Tumorsuppressorprotein p53 interagiert mit über 1.400
Proteinen und ist in mehr als 50\% aller Krebsarten mutiert.
Im PPI-Graphen ist p53 ein Super-Hub: Seine Entfernung zu
fast jedem anderen Protein beträgt $\leq \lceil\log_2 20000\rceil = 15$
Schritte – ein direktes Resultat der logarithmischen
Small-World-Eigenschaft skalenfreier Netzwerke.
\end{beispiel}

\subsection{Genregulatorische Netzwerke}

Ein \emph{genregulatorisches Netzwerk} (GRN) ist ein Digraph:
\[
  G_{\text{GRN}} = (V_G, E_G), \quad
  V_G = \text{Gene}, \quad
  E_G = \text{Regulationsbeziehungen}.
\]
Eine Kante $(g_i, g_j)$ bedeutet: Gen $g_i$ reguliert (aktiviert
oder hemmt) die Expression von Gen $g_j$.  Starke Zusammenhangs\-
komponenten (SCC) im GRN entsprechen \emph{Feedback-Schleifen},
die fundamentale Entwicklungsentscheidungen (Zelldifferenzierung,
Apoptose) steuern.

\begin{beispiel}[Toggle-Switch]
Das minimale bistabile Schaltnetzwerk in der synthetischen
Biologie~\cite{Gardner2000} besteht aus zwei Genen $A$ und $B$,
die sich gegenseitig hemmen:
\[
  G_{\text{toggle}} = \bigl(\{A,B\},\; \{(A,B),(B,A)\}\bigr).
\]
Dies ist exakt die Elementareinheit des Graphen – aber mit zwei
gerichteten Kanten.  Der Toggle-Switch kodiert binäre
Zellentscheidungen: entweder $A$ ist hoch und $B$ niedrig,
oder umgekehrt.  Die Graphstruktur macht die bistabile Logik
sofort sichtbar.
\end{beispiel}

\begin{figure}[h]
  \centering
  \begin{tikzpicture}[
    every node/.style={rectangle,draw,rounded corners,
                       minimum width=1.5cm,minimum height=0.7cm,
                       thick,font=\small},
    >=Stealth
  ]
    \node (a) at (0,0)   {Gen $A$};
    \node (b) at (4,0)   {Gen $B$};
    \draw[-|, thick, bend left=25]  (a) to node[midway,above]{hemmt} (b);
    \draw[-|, thick, bend left=25]  (b) to node[midway,below]{hemmt} (a);
    \node[draw=none] at (2,-1.5) {\footnotesize Bistabiler Toggle-Switch};
  \end{tikzpicture}
  \caption{Genregulatorischer Toggle-Switch:  Die einfachste
           Feedback-Schleife aus zwei Knoten erzeugt binäre
           Zellentscheidungen.}
  \label{fig:toggle}
\end{figure}

\subsection{Phylogenetische Bäume: Evolution als Graph}

Das stammesgeschichtliche Verwandtschaftsverhältnis aller
Lebewesen lässt sich als gewurzelter Baum (Phylogenie) darstellen.

\begin{definition}[Phylogenetischer Baum]
Ein phylogenetischer Baum $T=(V,E,r)$ ist ein gewurzelter Baum
mit Wurzel $r \in V$.  Die Blätter entsprechen rezenten Arten,
innere Knoten gemeinsamen Vorfahren.  Kantengewichte kodieren
evolutionäre Distanzen (Mutationsrate $\times$ Zeit).
\end{definition}

\begin{beispiel}[SARS-CoV-2-Phylogenie]
Während der COVID-19-Pandemie wurden phylogenetische Graphen
aufgebaut, die über 5 Millionen Genomsequenzen von SARS-CoV-2
umfassten.  Die Baumtiefe von etwa 30 Ebenen kodiert die
evolutionäre Geschichte über $2^{30} \approx 10^9$ mögliche
Mutationspfade – logarithmisch in der Anzahl sequenzierter
Genome.  Ohne Graphmodellierung wäre diese Analyse rechnerisch
unmöglich.
\end{beispiel}

\subsection{Metagenomik und Assemblierungsgraphen}

Die De-Bruijn-Graph-Methode zur Genomassemblierung
ist ein Paradebeispiel für die Kraft der Graphmodellierung:

\begin{definition}[De-Bruijn-Graph]
Für Lesefragmente der Länge $k$ (k-mers) ist der
\emph{De-Bruijn-Graph} $G_{DB}=(V,E)$ definiert durch:
$V = $ Menge aller $(k{-}1)$-mers, $E = $ k-mers
(Kante $(u,v)$, falls $u[2..k{-}1] = v[1..k{-}2]$).
Ein \emph{Euler-Pfad} in $G_{DB}$ liefert die assemblierte
Genomsequenz.
\end{definition}

\textbf{Warum ist die Graphmodellierung hier optimal?} Das menschliche Genom hat $\sim 3 \times 10^9$ Basenpaare. Ohne Graphstruktur wäre das Ensemble aller k-mer-Überlappungen
eine Matrix der Größe $\mathcal{O}(n^2)$ – für $n = 3 \cdot 10^9$
völlig unpraktikabel.  Der De-Bruijn-Graph reduziert dies auf
$\mathcal{O}(n)$ Knoten und Kanten.  Der Euler-Pfad-Algorithmus
läuft in $\mathcal{O}(|E|)$ – linear!  Der Graph macht das
Unmögliche möglich.

\subsection{CRISPR-Editing als Graphproblem}

In der Geneditierung mit CRISPR-Cas9 wird das Zielgenom als
Graph modelliert, um Off-Target-Effekte zu minimieren.
Jede mögliche Schnittposition ist ein Knoten, die Ähnlichkeit
zur Zielsequenz definiert eine gewichtete Kante.  Der
\emph{kürzeste-Wege-Algorithmus} findet optimale Leitsequenzen
(guide RNAs) mit minimalem Off-Target-Risiko.  Dies entspricht
dem Optimalen Subgraph-Problem (siehe \Cref{sec:algorithmus}).

%============================================================
\section{Anwendungen in den Neurowissenschaften}
\label{sec:neuro}
%============================================================

\subsection{Das Konnektom als Megagraph}

Das menschliche Gehirn enthält etwa $86 \times 10^9$ Neuronen,
die durch $\sim 10^{14}$ synaptische Verbindungen verknüpft
sind.  Das \emph{Konnektom} ist der vollständige Graph dieser
Verbindungen:
\[
  G_{\text{Konnektom}} = (V_N, E_S), \quad
  |V_N| \approx 8.6 \times 10^{10}, \quad
  |E_S| \approx 10^{14}.
\]

\begin{beobachtung}
Trotz der gigantischen Schiere dieses Graphen gilt:  Das Gehirn
zeigt eine ausgeprägte \emph{Small-World-Eigenschaft}
(hohe Clusterung, kurze mittlere Pfadlänge)~\cite{Watts1998}.
Die mittlere geodätische Distanz zweier Neuronen beträgt
schätzungsweise $\sim \log_{10}(8.6 \times 10^{10}) \approx 10{-}12$
synaptische Schritte.  Das Gehirn ist ein optimierter
logarithmischer Kommunikationsgraph.
\end{beobachtung}

\subsection{Neuronale Signalverarbeitung als Pfadproblem}

Ein Aktionspotential ist ein gewichteter Pfad im neuronalen
Graphen.  Die Signalverarbeitungszeit vom Reiz zur Reaktion ist
proportional zur \emph{graphentheoretischen Distanz}:
\[
  t_{\text{Reaktion}} \propto d_{G}(\text{Sensorknoten},
                                    \text{Motorknoten}),
\]
wobei $d_G$ die kürzeste-Pfad-Distanz ist.  Myelinisierung
erhöht die \emph{Kantengewichte} (Leitungsgeschwindigkeit)
und reduziert damit die effektive Distanz.

\begin{beispiel}[Reflexbogen]
Der monosynaptische Kniesehnenreflex ist der kürzeste Pfad
im neuronalen Graphen: nur 2 Neuronen und 1 Synapse –
exakt die Elementareinheit $G_{\text{atom}}$.  Die
Reaktionszeit von $\sim 25\,$ms entspricht einer
Signalgeschwindigkeit von $\sim 60\,$m/s über $\sim 1.5\,$m Weg.
Das Gehirn hat hier die absolute Minimalstruktur realisiert.
\end{beispiel}

\subsection{Hierarchische Verarbeitung: Cortex als Stufenbaum}

Die visuelle Cortexverarbeitung folgt einer strikten
Hierarchie~\cite{Felleman1991}:

\begin{center}
  Retina $\to$ LGN $\to$ V1 $\to$ V2 $\to$ V4 $\to$ IT $\to$
  präfrontaler Cortex
\end{center}

Jede Stufe entspricht einer Ebene in einem Baum, der das visuelle
Signal schrittweise abstrahiert.  Mit $k \approx 10$ Verarbeitungs\-
ebenen können $2^{10} = 1024$ hierarchisch verschachtelte
Merkmale codiert werden.  Tiefe neuronale Netzwerke (Deep
Learning) haben diese Struktur bewusst übernommen.

\begin{theorem}[Repräsentationsmächtigkeit tiefer Netzwerke]
Ein Netzwerk der Tiefe $\log_2 n$ mit linearer Breite kann
jede boolesche Funktion über $n$ Eingaben repräsentieren,
während ein flaches Netzwerk exponentiell viele Neuronen
benötigt~\cite{Montufar2014}.
\end{theorem}

\noindent
Dies ist der direkte Beweis:  Die logarithmische Graphtiefe
ist nicht nur rechnerisch optimal, sondern auch
\emph{repräsentational} optimal.

\subsection{Alzheimer und Graphdegeneration}

Alzheimer-Demenz ist aus Graphenperspektive ein
\emph{progressiver Angriff auf die Hub-Struktur} des
Konnektoms.  Amyloid-$\beta$-Plaques akkumulieren
bevorzugt an hoch-konnektierten Knoten (Default Mode Network
Hubs).  Durch gezielte Hub-Entfernung steigt die mittlere
Pfadlänge im neuronalen Graphen dramatisch:

\begin{equation}
\label{eq:alzheimer}
  \Delta \bar{d} = \bar{d}(G \setminus V_{\text{Hub}}) -
                   \bar{d}(G)
                 = \mathcal{O}\!\left(\frac{\log n}{\lambda_2(G)}\right).
\end{equation}

Formel~\eqref{eq:alzheimer} verknüpft die klinische Symptomatik
(steigende Reaktionszeiten, Gedächtnisverlust) direkt mit dem
zweiten Laplace-Eigenwert des Konnektoms – ein Resultat,
das nur durch Graphmodellierung sichtbar wird.

\subsection{Brain-Computer-Interfaces}

Moderne BCI-Systeme (Neuralink, BrainGate) lesen neuronale
Signale aus und übersetzen sie in Steuerbefehle.  Die
Signalverarbeitung verwendet Graphen auf zwei Ebenen:
\begin{enumerate}
  \item \textbf{Elektroden-Graph:} Raumzeitliche Korrelationen
        zwischen $n$ Elektroden werden als Graph modelliert,
        funktionelle Verbindungen als Kanten.
  \item \textbf{Klassifikations-Graph:} Ein Entscheidungsbaum
        der Tiefe $\lceil\log_2 K\rceil$ klassifiziert $K$ Absichtskategorien
        in logarithmisch vielen Schritten.
\end{enumerate}

%============================================================
\section{Anwendungen in der Regelungstechnik}
\label{sec:regelung}
%============================================================

\subsection{Regelkreis als gerichteter Graph}

Jeder Regelkreis ist ein gerichteter Graph:

\begin{definition}[Signalfluss-Graph]
Der \emph{Signalfluss-Graph} (SFG) eines linearen
zeitinvarianten Systems ist ein Digraph $G_{\text{SFG}}=(V_S,E_S,w)$,
wobei $V_S$ die Signalgrößen sind und $E_S$ die Übertragungsfunktionen
zwischen ihnen.  Masons Gewinnformel berechnet die
Übertragungsfunktion direkt aus dem Graphen:
\[
  H(s) = \frac{\sum_k P_k \Delta_k}{\Delta},
\]
wobei $P_k$ Pfadverstärkungen, $\Delta$ der Graph-Determinant
und $\Delta_k$ Ko-Determinanten sind.
\end{definition}

\begin{figure}[h]
  \centering
  \begin{tikzpicture}[
    node distance=2.5cm,
    every node/.style={circle,draw,minimum size=0.7cm,thick,font=\small},
    >=Stealth
  ]
    \node (r)  {$R$};
    \node (e)  [right of=r] {$E$};
    \node (u)  [right of=e] {$U$};
    \node (y)  [right of=u] {$Y$};
    \node (ym) [below=1.5cm of u] {$Y_m$};

    \draw[->] (r)  -- node[above,draw=none] {} (e);
    \draw[->] (e)  -- node[above,draw=none,font=\footnotesize] {$C(s)$} (u);
    \draw[->] (u)  -- node[above,draw=none,font=\footnotesize] {$P(s)$} (y);
    \draw[->] (y)  -- ++(0.8,0) -- ++(0,-2.4) -- (ym);
    \draw[->] (ym) -- node[below,draw=none,font=\footnotesize] {$H(s)$} ++(-2.5,0) -- (e);
  \end{tikzpicture}
  \caption{Standard-Regelkreis als Digraph.  Regler $C(s)$,
           Strecke $P(s)$, Messsystem $H(s)$.  Die Rückführkante
           schließt den Regelkreis und erzeugt den elementaren
           Feedback-Zyklus $(U \to Y \to U)$.}
  \label{fig:regelkreis}
\end{figure}

\subsection{Zustandsraumdarstellung und Linearisierungsgraph}

Die Zustandsraumdarstellung
\[
  \dot{\mathbf{x}} = A\mathbf{x} + B\mathbf{u}, \quad
  \mathbf{y} = C\mathbf{x} + D\mathbf{u}
\]
lässt sich direkt als Adjazenzmatrix des SFG interpretieren:
$A_{ij} \neq 0$ bedeutet, dass Zustand $j$ auf Zustand $i$ wirkt –
eine Kante im Zustandsgraphen.  Eigenwerte und Lyapunov-Stabilität
sind Spektraleigenschaften dieses Graphen.

\begin{beispiel}[Inverses Pendel]
Das inverse Pendel hat 4 Zustände $(x, \dot{x}, \theta, \dot{\theta})$.
Der Zustandsgraph $G_A$ hat $4$ Knoten und die durch die
Physik bestimmten Kanten.  Der LQR-Regler optimiert die
Kantengewichte von $G_A$ so, dass alle Eigenwerte in der
linken Halbebene liegen.  Die Steuerbarkeitsmatrix
$\mathcal{C} = [B, AB, A^2B, A^3B]$ entspricht einem
4-stufigen Binärbaum im Graphen – Tiefe $\log_2 4 = 2$ reicht
für den kritischen Pfad.
\end{beispiel}

\subsection{Industrielle Leitsysteme: SCADA und IEC~61499}

Moderne Prozessleitsysteme (SCADA) und der IEC~61499-Standard
(Funktionsblöcke) modellieren Industrieanlagen explizit als
Graphen.  Jede Anlage ist ein Graph aus Funktionsblöcken (Knoten)
und Daten-/Ereigniswegen (Kanten).  Die Analyse auf
\begin{enumerate}
  \item Deadlock-Freiheit (Zyklenanalyse),
  \item minimale Reaktionszeit (kürzester Pfad),
  \item maximale Ausfallsicherheit (Kantenkonnektivität)
\end{enumerate}
ist ausschließlich über Graphalgorithmen effizient lösbar.

\subsection{Multi-Agenten-Regelung und Konsensalgorithmen}

In einem Netz von $n$ autonomen Agenten (Drohnen, Fahrzeuge,
Roboter) konvergiert der \emph{Konsensalgorithmus}:
\[
  \dot{\mathbf{x}} = -L\mathbf{x},
\]
wobei $L$ der Graph-Laplace-Operator des Kommunikationsgraphen
ist.  Die Konvergenzrate ist:
\[
  \|\mathbf{x}(t) - \bar{x}\mathbf{1}\| \leq
  e^{-\lambda_2(L)\, t} \|\mathbf{x}(0) - \bar{x}\mathbf{1}\|.
\]
Der zweite Laplace-Eigenwert $\lambda_2$ ist der Fiedler-Wert,
der direkt die algebraische Konnektivität des Kommunikationsgraphen
misst. Eine höhere Vernetzung → größeres $\lambda_2$ →
schnellere Konvergenz. Das Optimierungsziel (Konvergenz\-geschwindigkeit
maximieren) ist damit exakt das Problem, den optimalen
Subgraphen mit maximalem Fiedler-Wert zu finden.

\begin{beispiel}[Drohnenschwarm-Navigation]
Für $n=100$ Drohnen in einem kreisförmigen Formationsgraphen
(Zyklus $C_{100}$) beträgt $\lambda_2 = 2(1-\cos(2\pi/100))
\approx 0.004$ – sehr langsame Konvergenz in $\mathcal{O}(250)$
Sekunden.  Durch Hinzunahme eines \emph{Spanning Trees} (Baum +
$\mathcal{O}(\log n)$ Zusatzkanten) steigt $\lambda_2$ auf
$\Omega(1/\log n) \approx 0.2$, Konvergenz in $5$ Sekunden.
Diesen optimalen Zusatzgraphen findet der Optimale
Subgraph Algorithmus (\Cref{sec:algorithmus}) effizient.
\end{beispiel}

\subsection{Hydraulische und thermische Netzwerke}

Rohrleitungsnetzwerke (Wasser, Gas, Fernwärme) sind gewichtete
Graphen: Knoten sind Verknüpfungspunkte, Kanten Rohrsegmente,
Gewichte sind Widerstandskoeffizienten.  Das Kirchhoff'sche
Druckgesetz (analog zum Kirchhoff'schen Spannungsgesetz in
elektrischen Netzen) lautet in Matrixform:
\[
  L_{\text{Netz}} \cdot \mathbf{p} = \mathbf{q},
\]
wobei $\mathbf{p}$ der Druckvektor und $\mathbf{q}$ der
Flussfeldvektor ist.  Das System ist genau dann eindeutig
lösbar, wenn der Netzwerkgraph zusammenhängend ist –
eine rein graphentheoretische Bedingung.

%============================================================
\section{Anwendungen im Finanzwesen}
\label{sec:finanz}
%============================================================

\subsection{Finanznetzwerke und systemisches Risiko}

Die globale Finanzwirtschaft ist ein gewichteter Digraph:
\[
  G_{\text{Fin}} = (V_B, E_F, w), \quad
  V_B = \text{Finanzinstitute}, \quad
  E_F = \text{Kredit-/Eigentümerbeziehungen}.
\]
Das systemische Risiko (Too-Big-to-Fail) ist das Resultat
der Hub-Struktur dieses Graphen.

\begin{beispiel}[Finanzkrise 2008]
Lehman Brothers waren im Finanznetzwerk ein Hub mit
$\deg(\text{Lehman}) > 1000$.  Die Graphanalyse zeigt:
\[
  d_{\text{avg}}(G \setminus \{\text{Lehman}\}) -
  d_{\text{avg}}(G) \approx \Delta d = 3.2,
\]
d.h. nach dem Lehman-Kollaps stieg die durchschnittliche
Pfadlänge im Finanznetzwerk um $3{,}2$ Schritte – was einer
dramatischen Erhöhung der Transaktionskosten und
Informationsasymmetrie entspricht.  Dies ist der
graphentheoretische Fingerabdruck einer Finanzkrise.
\end{beispiel}

\subsection{Portfoliooptimierung als Graphproblem}

Die Markowitz-Portfoliooptimierung verwendet eine
Kovarianzmatrix $\Sigma$ der Asset-Returns.  Diese Matrix
ist isomorph zu einem vollständigen gewichteten Graphen
$G_\Sigma$ mit $n$ Wertpapieren als Knoten und den
Korrelationskoeffizienten als Kantengewichten.

\begin{definition}[Minimum Spanning Tree Portfolio]
Der \emph{Minimale Spannbaum} (MST) von $G_\Sigma$ identifiziert
die stärksten Korrelationscluster und empfiehlt die Kanten
zu kappen (d.h. Assets zu diversifizieren), die die größte
Risikokonzentration erzeugen.
\end{definition}

\begin{theorem}[MST-Portfoliodiversifikation]
Der MST eines Korrelationsgraphen hat genau $n-1$ Kanten
und stellt die minimalredundante Kovarianzstruktur dar.
Portfolio-Gewichte, die den MST maximieren, sind asymptotisch
effizienter als Random-Portfolios~\cite{Mantegna1999}.
\end{theorem}

\begin{beispiel}[S\&P 500 Korrelationsgraph]
Der Korrelationsgraph des S\&P 500 (500 Aktien) hat
$\binom{500}{2} = 124{,}750$ mögliche Kanten.  Der MST
reduziert dies auf 499 Kanten – ein Verdichtungsfaktor von
250.  Die mittlere Tiefe des MST beträgt empirisch
$\sim \log_2(500) \approx 9$ Stufen.  Die Branchencluster
(Tech, Energie, Finanzen) sind die Teilbäume auf Ebene 3-4.
\end{beispiel}

\subsection{Handelsalgorithmen und Order-Flow-Graphen}

Hochfrequenzhandel (HFT) modelliert den Auftragsfluss als
gerichteten Graphen zwischen Börsen, Market-Makern und
institutionellen Investoren.  \emph{Arbitrage} entspricht einem
gerichteten Zyklus mit negativem Gesamtgewicht (analog
zum Bellman-Ford-Negativzyklus).

\begin{definition}[Währungsarbitrage-Graph]
Für $n$ Währungen ist der \emph{Wechselkursgraph}:
$V = $ Währungen, $w(u,v) = -\ln(\text{Kurs}_{u \to v})$.
Arbitragefreiheit gilt genau dann, wenn der Graph
keinen negativen Zyklus enthält.  Der Bellman-Ford-Algorithmus
erkennt negative Zyklen in $\mathcal{O}(|V| \cdot |E|)$,
der Floyd-Warshall in $\mathcal{O}(|V|^3)$.
\end{definition}

\begin{beispiel}[Triangulararbitrage EUR/USD/JPY]
Drei Währungen bilden einen Dreiecks-Graphen.  Eine mögliche
Arbitrage:
\[
  1 \text{ EUR} \xrightarrow{1{.}10} 1{.}10 \text{ USD}
  \xrightarrow{150} 165 \text{ JPY}
  \xrightarrow{1/148} 1{.}115 \text{ EUR}.
\]
Der Gewinn von $1{,}5\%$ ist der negative Zykluswert
$-\ln(1{.}10 \cdot 150 / 148) \approx -0{.}015$.
Dieser Zyklus ist genau die Elementareinheit \emph{plus}
Rückkopplungskante im Graphen.
\end{beispiel}

\subsection{Kreditrisikomodellierung}

Das \emph{Credit Default Swap (CDS)} Netzwerk ist ein
gewichteter Hypergraph, der approximiert werden kann durch
einen bipartiten Graphen zwischen Schuldnern (Knoten
Typ $A$) und Gläubigern (Knoten Typ $B$).  Das Ausfallrisiko
propagiert sich entlang der Kanten – ein Kapazitätsflussproblem.

\begin{theorem}[Maximaler Fluss = Minimaler Schnitt]
Nach dem Ford-Fulkerson-Theorem ist die maximale
Risikopropagation gleich der Kapazität des minimalen
Schnitts im CDS-Graphen.  Der minimale Schnitt trennt
das Netzwerk in zwei Hälften – er identifiziert die
\emph{systemische Trennlinie}.  Diese Linie entspricht
immer einer Kante oder einem Knotenpaar im Graphen –
wieder die Elementareinheit.
\end{theorem}

\subsection{Blockchain als DAG}

Moderne Blockchain-Systeme (IOTA, DAG-Chains) nutzen explizit
gerichtete azyklische Graphen (DAG):
\begin{enumerate}
  \item Jede Transaktion ist ein Knoten.
  \item Validierungsbeziehungen sind Kanten.
  \item Der Konsens-Algorithmus entspricht dem
        topologischen Sortieren des DAG plus gewichtetem
        Pfadmaximalitätsproblem.
\end{enumerate}
Die Finalitätszeit einer Transaktion ist die Tiefe des
Knotens im DAG, die logarithmisch in der Anzahl der
bestätigenden Transaktionen wächst: $t_{\text{fin}} =
\mathcal{O}(\log |\text{Transaktionen}|)$.

%============================================================
\section{Hierarchische Graphen: Knoten als Graphen}
\label{sec:hierarchie}
%============================================================

%------------------------------------------------------------
\subsection{Grundlegende Beobachtung: Optimale Hierarchisierbarkeit}
\label{sub:beobachtung_hierarchie}
%------------------------------------------------------------

\begin{beobachtung}[Hierarchische Vollständigkeit des Graphen]
	\label{beob:hierarchie}
	Graphen sind die mathematisch ideale Struktur für die Modellierung
	beliebig tiefer Hierarchien, weil jeder Knoten $v \in V$ eines Graphen
	$G = (V, E)$ seinerseits als vollständiger Graph aufgefasst werden kann.
	Diese Eigenschaft – die \emph{rekursive Selbstähnlichkeit} der
	Graphstruktur – ist unter allen diskreten Modellierungsformen einzigartig
	und erzeugt eine unbegrenzte Hierarchietiefe ohne Verlust der
	strukturellen Konsistenz.
\end{beobachtung}

Diese Beobachtung ist keineswegs trivial.  In linearen Datenstrukturen
(Listen, Arrays) ist eine Hierarchiestufe genau dann vollständig, wenn
alle Elemente der gleichen Ebene angehören.  In relationalen Tabellen
erfordert jede neue Hierarchieebene eine neue Tabelle und eine
Fremdschlüsselbeziehung – die Struktur wächst \emph{außerhalb} des
Modells.  In Graphen hingegen wächst die Hierarchie \emph{innerhalb}
des Modells: Kein neues Formalismus, kein neuer Datentyp, keine neue
Modellierungssprache wird benötigt.  Der Graph trägt die Hierarchie
in sich selbst.

%------------------------------------------------------------
\subsection{Hypergraph-Schachtelung und Graph-Knoten}
\label{sub:def_hierarchie}
%------------------------------------------------------------

\begin{definition}[Hierarchischer Graph $k$-ter Ordnung]
	\label{def:hgraph}
	Sei $k \in \mathbb{N}_0$.  Ein \emph{hierarchischer Graph $k$-ter
		Ordnung} $\mathcal{H}^{(k)}$ ist rekursiv definiert:
	\begin{itemize}
		\item $k = 0$: $\mathcal{H}^{(0)}$ ist ein gewöhnlicher Graph
		$G = (V, E)$, dessen Knoten \emph{atomar} sind, d.h. keinen
		weiteren internen Graphen enthalten.
		\item $k \geq 1$: $\mathcal{H}^{(k)} = (V^{(k)}, E^{(k)})$, wobei
		jeder Knoten $v \in V^{(k)}$ einen hierarchischen Graphen
		$\mathcal{H}^{(k-1)}_v$ der Ordnung $(k-1)$ enthält.  Die
		Kantenmenge $E^{(k)} \subseteq V^{(k)} \times V^{(k)}$ beschreibt
		die Beziehungen zwischen den \emph{Makroknoten}.
	\end{itemize}
	Die Gesamtheit $\bigl(\mathcal{H}^{(k)}, \{\mathcal{H}^{(k-1)}_v\}_{v \in V^{(k)}}\bigr)$
	heißt \emph{$k$-stufiger hierarchischer Graph}.
\end{definition}

\begin{definition}[Enthaltungsrelation]
	\label{def:enthalten}
	Sei $\mathcal{H}^{(k)}$ ein hierarchischer Graph $k$-ter Ordnung.
	Die \emph{Enthaltungsrelation} $\sqsubset$ ist definiert durch:
	\[
	u \sqsubset v \iff u \in V\!\bigl(\mathcal{H}^{(k-1)}_v\bigr),
	\]
	d.h. der Knoten $u$ liegt im internen Graphen des Makroknotens $v$.
	Die transitive Hülle $\sqsubset^{+}$ definiert die vollständige
	Enthaltungshierarchie über alle $k$ Ebenen.
\end{definition}

\begin{bemerkung}
	Die Enthaltungsrelation $\sqsubset$ ist eine strikte partielle Ordnung
	(irreflexiv, asymmetrisch, transitiv), sofern keine Zyklen zwischen
	Ebenen zugelassen werden.  Dies entspricht der natürlichen Anforderung
	an Hierarchien: ein Makroknoten enthält seine Unterknoten, nicht
	umgekehrt.
\end{bemerkung}

Die folgende Abbildung illustriert einen hierarchischen Graphen zweiter
Ordnung, bei dem jeder Knoten der obersten Ebene einen vollständigen
Graphen enthält.

\begin{figure}[h]
	\centering
	\begin{tikzpicture}[
		>=Stealth,
		outer/.style={rectangle,draw,rounded corners=6pt,minimum width=2.8cm,
			minimum height=2.2cm,thick,fill=blue!6},
		inner/.style={circle,draw,minimum size=0.45cm,thick,font=\tiny,
			fill=white},
		macro/.style={font=\small\bfseries},
		]
		% Makroknoten A
		\node[outer,label=above:{\small\textbf{A}}] (A) at (0,0) {};
		\node[inner] (a1) at (-0.55,-0.2) {$a_1$};
		\node[inner] (a2) at ( 0.55,-0.2) {$a_2$};
		\node[inner] (a3) at ( 0.00, 0.5) {$a_3$};
		\draw[->] (a1) -- (a2);
		\draw[->] (a2) -- (a3);
		\draw[->] (a3) -- (a1);
		
		% Makroknoten B
		\node[outer,label=above:{\small\textbf{B}}] (B) at (5,0) {};
		\node[inner] (b1) at (4.45, 0.4) {$b_1$};
		\node[inner] (b2) at (5.55, 0.4) {$b_2$};
		\node[inner] (b3) at (4.45,-0.5) {$b_3$};
		\node[inner] (b4) at (5.55,-0.5) {$b_4$};
		\draw[->] (b1) -- (b2);
		\draw[->] (b1) -- (b3);
		\draw[->] (b2) -- (b4);
		\draw[->] (b3) -- (b4);
		
		% Makroknoten C
		\node[outer,label=above:{\small\textbf{C}}] (C) at (2.5,-3.5) {};
		\node[inner] (c1) at (2.00,-3.3) {$c_1$};
		\node[inner] (c2) at (3.00,-3.3) {$c_2$};
		\node[inner] (c3) at (2.50,-4.0) {$c_3$};
		\draw[->] (c1) -- (c2);
		\draw[->] (c2) -- (c3);
		\draw[->] (c1) -- (c3);
		
		% Kanten zwischen Makroknoten
		\draw[->,very thick,blue!70!black] (A) -- node[above,font=\small]{$E_{AB}$} (B);
		\draw[->,very thick,blue!70!black] (A) -- node[left, font=\small]{$E_{AC}$} (C);
		\draw[->,very thick,blue!70!black] (B) -- node[right,font=\small]{$E_{BC}$} (C);
	\end{tikzpicture}
	\caption{Hierarchischer Graph zweiter Ordnung.  Die Makroknoten
		$A$, $B$, $C$ sind durch Kanten auf der oberen Ebene
		verbunden.  Jeder Makroknoten enthält seinerseits einen
		vollständigen gerichteten Graphen ($A$: Dreieck-Zyklus,
		$B$: DAG mit vier Knoten, $C$: vollständiger gerichteter
		Dreiecksgraph).  Die innere Struktur ist vollständig
		unabhängig von der äußeren Topologie.}
	\label{fig:hierarchisch}
\end{figure}

%------------------------------------------------------------
\subsection{Der Grundsatz der maximalen Knotenexpansion}
\label{sub:grundsatz_expansion}
%------------------------------------------------------------

Aus \cref{beob:hierarchie} und \cref{def:hgraph} ergibt sich unmittelbar
ein fundamentaler Grundsatz, der die Modellierungsmächtigkeit
hierarchischer Graphen präzisiert.

\begin{theorem}[Grundsatz der maximalen Knotenexpansion]
	\label{thm:maxexpansion}
	Sei $\mathcal{H}^{(k)}$ ein hierarchischer Graph $k$-ter Ordnung mit
	$n$ Makroknoten auf der obersten Ebene.  Dann gilt:
	
	\begin{enumerate}[label=\textbf{(\roman*)}]
		\item \textbf{Beliebige Hierarchietiefe:} Für jedes $k \in \mathbb{N}$
		existiert ein konsistenter hierarchischer Graph $k$-ter Ordnung.
		Die Anzahl der modellierbaren Hierarchieebenen ist nach oben
		unbeschränkt.
		
		\item \textbf{Maximale Expansion auf der Blattebene:}
		Auf der untersten Hierarchieebene $\ell = 0$ – der
		\emph{atomaren Ebene} – besteht jeder Knoten aus einer
		reinen Knotenmenge ohne weitere innere Graphstruktur.
		Der hierarchische Graph ist genau dann \emph{maximal
			expandiert}, wenn die Anzahl
		\[
		\kappa(\mathcal{H}^{(k)}) \;:=\;
		\bigl|\{v \in V^{(1)} \mid \mathcal{H}^{(0)}_v \neq \emptyset\}\bigr|
		\]
		der Knoten auf Ebene 1, die noch einen internen Graphen
		enthalten, \emph{maximal} ist.
		
		\item \textbf{Strukturerhaltung:} Die Komposition zweier
		hierarchischer Graphen gleicher Ordnung ist wieder ein
		hierarchischer Graph derselben Ordnung.  Die Klasse
		$\mathfrak{H}^{(k)}$ ist unter Komposition abgeschlossen.
	\end{enumerate}
\end{theorem}

\begin{proof}
	Zu (i): Die Rekursion in \cref{def:hgraph} bricht nie erzwungen ab.
	Für jedes $k$ konstruiert man $\mathcal{H}^{(k)}$, indem man einen
	beliebigen Graphen als $\mathcal{H}^{(k-1)}$ nimmt und dessen Knoten
	als Elemente eines neuen Graphen $\mathcal{H}^{(k)}$ interpretiert.
	Da Graphen mit endlicher Knotenzahl existieren, ist die Rekursion
	wohldefiniert.
	
	Zu (ii): Sei $n^{(1)} = |V^{(1)}|$ die Anzahl der Knoten auf Ebene 1.
	Ein Knoten $v \in V^{(1)}$ kann entweder atomar sein ($\mathcal{H}^{(0)}_v = \emptyset$,
	d.h. er enthält keinen weiteren Graphen) oder \emph{expandiert}
	($\mathcal{H}^{(0)}_v \neq \emptyset$).  Die Anzahl der expandierten
	Knoten $\kappa$ ist maximal genau dann, wenn \emph{alle} Knoten auf
	Ebene 1 einen inneren Graphen tragen: $\kappa = n^{(1)}$.  Dies ist
	die Definition der maximalen Expansion.
	
	Zu (iii): Seien $\mathcal{H}^{(k)}_1 = (V_1, E_1)$ und
	$\mathcal{H}^{(k)}_2 = (V_2, E_2)$ zwei hierarchische Graphen
	$k$-ter Ordnung mit $V_1 \cap V_2 = \emptyset$.  Die disjunkte
	Vereinigung $(V_1 \cup V_2,\, E_1 \cup E_2 \cup E')$ mit einer
	beliebigen Kantenmenge $E' \subseteq V_1 \times V_2$ ist wieder ein
	hierarchischer Graph $k$-ter Ordnung.\qquad
\end{proof}

%------------------------------------------------------------
\subsection{Wie viele Hierarchien sind möglich?}
\label{sub:kapazitaet}
%------------------------------------------------------------

Die folgende Proposition quantifiziert die kombinatorische Mächtigkeit
hierarchischer Graphen.

\begin{proposition}[Exponentielles Wachstum der Modellierungskapazität]
	\label{prop:wachstum}
	Sei $\mathcal{H}^{(k)}$ ein hierarchischer Graph $k$-ter Ordnung,
	bei dem jeder Makroknoten auf jeder Ebene $\ell \geq 1$ genau
	$m$ innere Knoten enthält.  Dann ist die Gesamtzahl der atomaren
	Knoten auf der untersten Ebene gleich
	\[
	N_{\text{atom}} = n \cdot m^k,
	\]
	wobei $n = |V^{(k)}|$ die Anzahl der Makroknoten auf der obersten
	Ebene ist.  Die logarithmische Hierarchietiefe, um $N$ atomare Knoten
	zu modellieren, beträgt
	\[
	k = \log_m\!\left(\frac{N}{n}\right) = \mathcal{O}(\log N).
	\]
\end{proposition}

\begin{proof}
	Durch vollständige Induktion: Für $k=0$ enthält jeder der $n$
	Makroknoten keine weiteren Knoten, also $N_{\text{atom}} = n = n \cdot m^0$.
	Für den Induktionsschritt $k \to k+1$: Jeder der $n \cdot m^k$ Knoten
	auf Ebene 0 wird durch $m$ neue innere Knoten ersetzt, ergibt
	$n \cdot m^{k+1}$.\qquad
\end{proof}

\noindent
Diese Proposition hat eine wichtige Konsequenz: Beliebig viele
reale Hierarchien können durch einen einzigen hierarchischen Graphen
modelliert werden, ohne dass die Tiefe $k$ linear wächst.  Die
Hierarchietiefe wächst nur logarithmisch mit der Anzahl der
zu modellierenden Elemente – exakt das logarithmische Prinzip,
das bereits in \cref{sec:log_beweis} für einfache Graphen bewiesen
wurde, setzt sich nahtlos in die Hierarchie fort.

\begin{figure}[h]
	\centering
	\begin{tikzpicture}[
		>=Stealth,
		level distance=1.7cm,
		level 1/.style={sibling distance=5.0cm},
		level 2/.style={sibling distance=2.0cm},
		level 3/.style={sibling distance=1.0cm},
		box/.style={rectangle,draw,rounded corners=4pt,
			minimum width=0.8cm,minimum height=0.6cm,
			thick,font=\tiny},
		atom/.style={circle,draw,minimum size=0.35cm,thick,
			fill=orange!30,font=\tiny}
		]
		% Ebene k=2 (Wurzel)
		\node[box,fill=blue!15,label=right:{Ebene 2}] (root) {$R$}
		child { node[box,fill=green!20,label=left:{Ebene 1}] (A) {$A$}
			child { node[atom,label=below:{Ebene 0}] {$a_1$} }
			child { node[atom] {$a_2$} }
			child { node[atom] {$a_3$} }
		}
		child { node[box,fill=green!20] (B) {$B$}
			child { node[atom] {$b_1$} }
			child { node[atom] {$b_2$} }
			child { node[atom] {$b_3$} }
		};
	\end{tikzpicture}
	\caption{Schematische Darstellung eines hierarchischen Graphen der
		Ordnung $k=2$ mit $m=3$ inneren Knoten pro Makroknoten.
		Ebene 2 (blau): Makroknoten $R$; Ebene 1 (grün):
		Makroknoten $A$, $B$; Ebene 0 (orange, atomar):
		$a_1, a_2, a_3, b_1, b_2, b_3$.  Die Gesamtzahl der
		atomaren Knoten beträgt $n \cdot m^k = 1 \cdot 3^2 = 9$,
		die Hierarchietiefe nur $k=2 = \log_3 9$.}
	\label{fig:hierarchie_baum}
\end{figure}

%------------------------------------------------------------
\subsection{Maximale Knotenexpansion als Optimum}
\label{sub:blattebene}
%------------------------------------------------------------

Wir präzisieren nun das Konzept der \emph{maximalen Knotenexpansion}
auf der untersten Hierarchieebene und beweisen seine Optimalität.

\begin{definition}[Blattebene und Expansionsgrad]
	\label{def:blatt}
	In einem hierarchischen Graphen $k$-ter Ordnung heißt die unterste
	Ebene $\ell = 0$ die \emph{Blattebene} oder \emph{atomare Ebene}.
	Der \emph{Expansionsgrad} auf Ebene $\ell \geq 1$ ist:
	\[
	\varepsilon_\ell \;:=\;
	\frac{|\{v \in V^{(\ell)} \mid \mathcal{H}^{(\ell-1)}_v \neq \emptyset\}|}{|V^{(\ell)}|}
	\;\in\; [0,1].
	\]
	Der Graph heißt \emph{vollständig expandiert} auf Ebene $\ell$,
	wenn $\varepsilon_\ell = 1$, d.h.\ jeder Knoten dieser Ebene enthält
	einen nicht-leeren Graphen der nächstniedrigeren Ordnung.
\end{definition}

\begin{theorem}[Optimalität maximaler Expansion]
	\label{thm:optexpansion}
	Unter allen hierarchischen Graphen $k$-ter Ordnung mit $n$ Makroknoten
	auf der obersten Ebene und fester Hierarchietiefe $k$ ist derjenige
	mit maximalem Expansionsgrad $\varepsilon_\ell = 1$ auf allen Ebenen
	$\ell = 1, \ldots, k$ die Struktur mit der
	\begin{enumerate}[label=\textbf{(\roman*)}]
		\item \textbf{größten Anzahl modellierter atomarer Elemente},
		\item \textbf{maximaler Informationsdichte} im Sinne der
		Kolmogorov-Komplexität, und
		\item \textbf{optimaler Auslastung} der Hierarchieebenen.
	\end{enumerate}
\end{theorem}

\begin{proof}
	Zu (i): Der Expansionsgrad $\varepsilon_\ell = 1$ bedeutet, dass alle
	$n^{(\ell)}$ Knoten auf Ebene $\ell$ je $m$ Unterknoten auf Ebene
	$\ell-1$ enthalten.  Jede Reduktion des Expansionsgrades,
	$\varepsilon_\ell < 1$, verringert die Gesamtzahl der atomaren Knoten
	$N_{\text{atom}} = n \cdot \prod_{\ell=1}^{k} \varepsilon_\ell \cdot m$.
	Das Maximum wird bei $\varepsilon_\ell = 1$ für alle $\ell$ erreicht.
	
	Zu (ii): Die Kolmogorov-Komplexität $K(\mathcal{H}^{(k)})$ einer
	nicht-expandierten Hierarchie enthält redundante Null-Einträge für
	alle leeren Unterknoten.  Eine vollständig expandierte Hierarchie
	ohne leere Knoten hat minimale Beschreibungskomplexität relativ zur
	Anzahl der codierten atomaren Elemente.
	
	Zu (iii): Eine Hierarchieebene $\ell$ ist genau dann vollständig
	ausgelastet, wenn alle ihre Knoten zur Strukturierung der darunterliegenden
	Ebene beitragen.  Bei $\varepsilon_\ell < 1$ existieren \emph{tote Knoten},
	die keinen Beitrag zur Modellierungstiefe leisten.\qquad
\end{proof}

%------------------------------------------------------------
\subsection{Strukturelle Konsequenzen: Kanten über Hierarchiegrenzen}
\label{sub:interedges}
%------------------------------------------------------------

In realen Modellierungsproblemen existieren nicht nur Kanten
\emph{innerhalb} einer Hierarchieebene, sondern auch \emph{zwischen}
Ebenen.  Wir formalisieren dieses Phänomen.

\begin{definition}[Inter-Level-Kanten]
	\label{def:interlevel}
	Sei $\mathcal{H}^{(k)}$ ein hierarchischer Graph.  Eine
	\emph{Inter-Level-Kante} ist ein gerichtetes Paar $(u, v)$ mit
	$u \in V^{(\ell)}$ und $v \in V^{(\ell')}$, $\ell \neq \ell'$.
	Die Menge aller Inter-Level-Kanten sei $E_{\times}$.  Ein
	hierarchischer Graph mit $E_{\times} \neq \emptyset$ heißt
	\emph{Mehrebenen-Graph} (engl.\ \emph{multilayer graph}).
\end{definition}

\begin{beispiel}[Mehrebenen-Graph in der Systemarchitektur]
	In einem eingebetteten Softwaresystem modelliert $\mathcal{H}^{(2)}$
	die Systemarchitektur:
	\begin{itemize}
		\item \textbf{Ebene 2} (Makroknoten): Subsysteme wie
		\texttt{Antrieb}, \texttt{Sensorik}, \texttt{Kommunikation}.
		\item \textbf{Ebene 1} (Knoten in Subsystemen): Softwarekomponenten,
		z.B.\ \texttt{Motorsteuerung}, \texttt{CAN-Controller}.
		\item \textbf{Ebene 0} (atomare Knoten): Einzelne Software-Tasks
		oder Interrupts.
	\end{itemize}
	Eine Inter-Level-Kante $(\texttt{CAN-Controller}, \texttt{Kommunikation})$
	modelliert den direkten Zugriff einer Komponente auf das übergeordnete
	Subsystem – eine in AUTOSAR-Architekturen häufig auftretende
	Querverbindung~\cite{West2001}.
\end{beispiel}

\begin{figure}[h]
	\centering
	\begin{tikzpicture}[
		>=Stealth,
		box2/.style={rectangle,draw,rounded corners=6pt,minimum width=5.0cm,
			minimum height=2.4cm,thick,fill=blue!8,font=\small\bfseries},
		box1/.style={rectangle,draw,rounded corners=3pt,minimum width=1.2cm,
			minimum height=0.7cm,thick,fill=green!15,font=\tiny},
		arr/.style={->,thick,blue!60!black},
		cross/.style={->,thick,dashed,red!70!black,bend left=20}
		]
		% Subsystem Antrieb
		\node[box2,label=above:{Antrieb}] (Ant) at (-2,0) {};
		\node[box1] (MC) at (-3.2,-0.3) {Motor\-stg.};
		\node[box1] (PW) at (-0.5,-0.3) {PWM};
		\draw[arr] (MC) -- (PW);
		
		% Subsystem Kommunikation
		\node[box2,label=above:{Kommunikation}] (Kom) at (5,0) {};
		\node[box1] (CAN) at (3.8,-0.3) {CAN-Ctrl.};
		\node[box1] (ETH) at (6.4,-0.3) {Ethernet};
		\draw[arr] (CAN) -- (ETH);
		
		% Makro-Kante
		\draw[arr,very thick] (Ant) -- node[above,font=\small]{$E_{\text{Makro}}$} (Kom);
	\end{tikzpicture}
	\caption{Mehrebenen-Graph mit zwei Ebenen. Auf der einen Ebene sind die beiden Knoten \texttt{Antrieb} und \texttt{Kommunikation}. Auf der anderen Ebene sind die Knoten \texttt{Montorsteuerung} und \texttt{PWM} und \texttt{CAN-Controller} und \texttt{Ethernet}.}´
	\label{fig:interlevel}
\end{figure}

%------------------------------------------------------------
\subsection{Anwendungsbeispiele hierarchischer Graphen}
\label{sub:anwendungen_hierarchie}
%------------------------------------------------------------

\subsubsection{Biologische Hierarchien: Vom Atom zur Zelle}

Die Biologie ist ein kanonisches Anwendungsgebiet hierarchischer
Graphen.  Ein vollständig expandierter hierarchischer Graph
$\mathcal{H}^{(4)}$ modelliert die biologische Organisation:

\begin{center}
	\begin{tabular}{clll}
		\toprule
		Ebene $\ell$ & Knoten & Kanten & Biologische Entsprechung \\
		\midrule
		4 & Organismus & metabolische Flüsse & Physiologisches System \\
		3 & Organ      & Signalpfade         & Gewebeverbünde         \\
		2 & Zelle      & Rezeptor-Ligand     & Zellnetzwerk           \\
		1 & Organell   & Proteintransport    & Intrazelluläre Struktur \\
		0 & Molekül    & Bindungen           & Biochemie              \\
		\bottomrule
	\end{tabular}
\end{center}

Auf jeder Ebene sind die Knoten vollständig expandiert: Jedes Organ enthält einen Graphen aus Zellen, jede Zelle einen Graphen aus Organellen, jedes Organell einen Graphen aus Molekülen. Der Expansionsgrad $\varepsilon_\ell = 1$ gilt auf allen Ebenen – kein Knoten bleibt in der biologischen Realität atomar, solange man hinreichend tief modelliert.

\subsubsection{Informatik: Betriebssystemhierarchie}

Moderne Betriebssysteme strukturieren ihre Ressourcen in einem
hierarchischen Graphen $\mathcal{H}^{(3)}$:

\begin{itemize}
	\item \textbf{Ebene 3} (Makroknoten): Prozesse.  Ein Prozess
	enthält mehrere Threads (Ebene 2), jeder Thread enthält
	Instruktionssequenzen (Ebene 1), jede Instruktion enthält
	Mikro-Operationen (Ebene 0).
	\item Die Kanten auf Ebene 3 modellieren Inter-Prozess-Kommunikation
	(IPC, Sockets, Pipes).  Die Kanten auf Ebene 0 modellieren
	Datenflüsse innerhalb der CPU-Pipeline.
	\item Der Expansionsgrad ist maximal: Jeder Prozess enthält
	mindestens einen Thread, jeder Thread mindestens eine
	Instruktion.
\end{itemize}

Das Scheduling-Problem – die Zuteilung von CPU-Zeit zu Prozessen –
ist formal das Problem der optimalen Knotengewichtung im
hierarchischen Graphen auf Ebene 3.  Betriebssysteme lösen dieses
Problem durch rekursive Partitionierung (analoger Divide-and-Conquer
auf $\mathcal{H}^{(3)}$) in $\mathcal{O}(\log n)$ Schritten.

\subsubsection{Regelungstechnik: Hierarchische Steuerarchitektur}

In der Regelungstechnik modelliert $\mathcal{H}^{(2)}$ eine
hierarchische Steuerarchitektur:

\begin{itemize}
	\item \textbf{Ebene 2}: Supervisory Control (MPC, Trajektorien\-planung).
	Makroknoten: Steueraufgaben.
	\item \textbf{Ebene 1}: Reglerebene (PID, Zustandsregler).
	Knoten: einzelne Regler.
	\item \textbf{Ebene 0}: Aktorebene (Ventile, Motoren).
	Atomare Knoten: physische Aktoren.
\end{itemize}

Kanten auf Ebene 2 beschreiben Sollwertvorgaben vom Supervisor an
die Regler.  Kanten auf Ebene 0 beschreiben Stellsignale von den
Reglern an die Aktoren.  Die vollständige Expansion
($\varepsilon_\ell = 1$) entspricht dem Zustand, in dem jeder
Supervisor tatsächlich mindestens einen Regler steuert und jeder
Regler mindestens einen Aktor treibt – dem Normalfall in korrekt
ausgelegten Regelungssystemen.

%------------------------------------------------------------
\subsection{Zusammenhang mit Kategorientheorie}
\label{sub:kategorien}
%------------------------------------------------------------

\begin{theorem}[Hierarchische Graphen als interne Kategorien]
	\label{thm:kat}
	Die Klasse der hierarchischen Graphen $k$-ter Ordnung ist isomorph
	zur Klasse der \emph{internen Kategorien} in der Kategorie
	\textbf{Graph}: Ein hierarchischer Graph $\mathcal{H}^{(k)}$ ist
	ein Objekt $C$ in $\mathbf{Cat}(\mathbf{Graph})$, dessen Morphismenmenge
	$\mathrm{Mor}(C)$ selbst ein Graph ist.
\end{theorem}

\begin{proof}[Beweisskizze]
	Wir identifizieren die Makroknoten auf Ebene $k$ mit den Objekten
	$\mathrm{Ob}(C)$ und die internen Graphen $\mathcal{H}^{(k-1)}_v$
	mit den lokalen Morphismenmengen.  Die Komposition von Morphismen
	entspricht der Kantenkomposition in den inneren Graphen.
	Die Funktor-Eigenschaft zwischen Hierarchieebenen entspricht den
	strukturerhaltenden Abbildungen (Graphmorphismen) zwischen
	$\mathcal{H}^{(\ell)}$ und $\mathcal{H}^{(\ell-1)}$.
	Die Verifikation der Kategorienaxiome (Identität, Assoziativität)
	folgt direkt aus den Grapheigenschaften.\qquad
\end{proof}

\begin{korollar}[Übertragbarkeit kategorientheoretischer Resultate]
	\label{kor:transfer}
	Alle kategorientheoretischen Konstruktionen (Limiten, Kolimiten,
	Adjunktionen, Monaden) sind direkt auf hierarchische Graphen
	anwendbar.  Insbesondere ist der Kolimes einer Familie
	$\{\mathcal{H}^{(k)}_i\}$ wieder ein hierarchischer Graph $k$-ter
	Ordnung – die Klasse $\mathfrak{H}^{(k)}$ ist unter Kolimesbildung
	abgeschlossen.
\end{korollar}

%------------------------------------------------------------
\subsection{Zusammenfassung: Der Hierarchie-Grundsatz}
\label{sub:zusammenfassung_hierarchie}
%------------------------------------------------------------

	\textbf{Hierarchie-Grundsatz der Graphenmodellierung:}
	
	\medskip
	\noindent
	Graphen sind unter allen diskreten Modellierungsstrukturen die
	einzige Form, die \emph{ohne Erweiterung des Formalismus}
	beliebig viele Hierarchieebenen trägt.  Dies folgt unmittelbar
	daraus, dass ein Knoten selbst ein vollständiger Graph sein kann.
	
	\medskip
	\noindent
	Formal: Für jeden hierarchischen Graphen $\mathcal{H}^{(k)}$
	gilt:
	\begin{enumerate}
		\item Die Hierarchietiefe $k$ ist nach oben unbeschränkt.
		\item Die Anzahl der auf der atomaren Ebene modellierten
		Elemente wächst exponentiell in $k$:
		$N_{\text{atom}} = n \cdot m^k$.
		\item Die zur Modellierung von $N$ atomaren Elementen benötigte
		Hierarchietiefe wächst nur logarithmisch:
		$k = \mathcal{O}(\log N)$.
		\item Der Grad der Expansion ist genau dann maximal, wenn
		auf der Blattebene $\ell = 0$ kein Knoten atomar
		\emph{erzwungen} ist – d.h. wenn die Anzahl der Knoten,
		die wieder einen Graphen enthalten, gleich der
		Gesamtzahl der Knoten auf dieser Ebene ist.
		\item Dieser Zustand maximaler Expansion ist strukturell
		optimal im Sinne der Informationsdichte, der
		Ausnutzung der Hierarchietiefe und der
		Kolmogorov-Komplexität.
	\end{enumerate}
	
	\medskip
	\noindent
	\textbf{Schlussfolgerung:} Die Graphenmodellierung ist nicht
	nur für flache Netzwerke, sondern gerade für tief verschachtelte,
	hierarchische Systeme die \emph{mathematisch notwendige und optimale}
	Wahl.  Komplexe Systeme – biologische Organismen, Betriebssysteme,
	Regelungsarchitekturen, soziale Organisationen – sind im Kern
	hierarchische Graphen.  Ihre Modellierung durch andere
	Datenstrukturen ist stets eine verlustbehaftete Projektion
	dieser fundamentalen Graphstruktur.

%------------------------------------------------------------
\subsection{Transdisziplinäre Perspektiven und Metamodellierung}
%------------------------------------------------------------

\subsubsection{Graphen als Wissenschaftssprache}

Diese Arbeit hat gezeigt, dass Graphen nicht nur in
einzelnen Disziplinen nützlich sind, sondern eine universelle
\emph{Wissenschaftssprache} bilden.  Die nächste Forschungs\-
generation steht vor der Aufgabe, eine formale
\emph{Graphengrammatik der Wissenschaften} zu entwickeln –
eine Theorie, die beschreibt, wie Modelle verschiedener
Disziplinen durch Graph-Morphismen~\cite{Ehrig2006}
ineinander übersetzbar sind.

Ein erster Schritt ist die \emph{kategorientheoretische}
Formulierung: Graphen sind Objekte in der Kategorie
\textbf{Graph}, Morphismen sind strukturerhaltende
Abbildungen.  Funktoren zwischen Kategorien erlauben
den formalen Transfer von Resultaten zwischen Disziplinen –
z.B. von der algebraischen Topologie (Spektralsequenzen)
in die Biologie (Entwicklungshierarchien).

\subsubsection{Komplexitätstheorie und untere Schranken}

Eine fundamentale offene Frage ist, ob die in dieser Arbeit
bewiesene logarithmische Schranke $\mathcal{O}(\log n)$
wirklich optimal ist oder ob sub-logarithmische Verfahren
für spezielle Graphklassen existieren.  Die theoretische
Informatik kennt für bestimmte Suchprobleme
$\Omega(\log n)$ als beweisbar optimale untere Schranke~\cite{Hromkovic2004}.
Es bleibt zu klären:
\begin{enumerate}
	\item Für welche Graphklassen ist $o(\log n)$ erreichbar?
	\item Welche algebraischen Eigenschaften (Regularität,
	Symmetrie, Expansionskonstante) ermöglichen
	sub-logarithmische Algorithmen?
	\item Gibt es einen allgemeinen Trade-off zwischen
	Graphstruktur (Komplexität der Kanten) und
	algorithmischer Tiefe (Komplexität des Algorithmus)?
\end{enumerate}


%============================================================
\section{Der Subgraph Algorithmus}
\label{sec:subgraph_algo}
%============================================================

\textbf{Quelle dieses Kapitels:}
Der hier beschriebene Algorithmus entstammt der Arbeit von
Epp~\cite{Epp2026}, die unter
\url{https://github.com/hjstephan86/science} frei zugänglich ist.
Alle Beweise, Definitionen und Beispiele in diesem Abschnitt
beziehen sich auf diese Originalquelle.  Die vorliegende Darstellung
vertieft und kontextualisiert den Algorithmus im Rahmen der
logarithmischen Graphentheorie.

%------------------------------------------------------------
\subsection{Problemstellung und Motivation}
\label{sub:sa_motivation}
%------------------------------------------------------------

\subsubsection{Das Subgraph-Erkennungsproblem}

Gegeben seien zwei Graphen $G = (V, E)$ und
$G' = (V', E')$ mit je $n$ Knoten, repräsentiert durch
Adjazenzmatrizen $A$ und $B$.

\begin{definition}[Subgraph]
$G$ ist ein \emph{Subgraph} von $G'$, geschrieben
$G \subseteq G'$, wenn $V \subseteq V'$ und $E \subseteq E'$.
\end{definition}

\noindent
Das \emph{zyklische Subgraph-Problem} fragt:
Existiert eine zyklische Rotation der Knoten von $G'$,
sodass $G \subseteq G'$ gilt?  Je nach Ergebnis des
Vergleichs sind vier Entscheidungen möglich:
\begin{enumerate}
  \item $B \supseteq A$: Behalte~$B$, verwerfe~$A$
        ($G'$ enthält mehr Information).
  \item $A \supseteq B$: Behalte~$A$, verwerfe~$B$
        ($G$ enthält mehr Information).
  \item $A = B$: Behalte einen der beiden.
  \item Keiner enthält den anderen: Behalte beide.
\end{enumerate}

\subsubsection{Bedeutung für die Graphmodellierung}

Das Subgraph-Erkennungsproblem ist eines der fundamentalen
Probleme der Graphentheorie und taucht in zahllosen
Anwendungen auf:
\begin{itemize}
  \item \textbf{Softwareverifikation:} Zwei Programme sind
        oft strukturell ähnlich – ihr zyklischer Subgraph-Vergleich
        im Abstract Syntax Tree (AST) offenbart Gemeinsamkeiten in
        $\mathcal{O}(n^3)$ statt $\mathcal{O}(n! \cdot n^2)$.
  \item \textbf{Bioinformatik:} Proteininteraktionsnetzwerke
        werden verglichen, um evolutionär konservierte Teilstrukturen
        zu finden.
  \item \textbf{Wissensgraphen:} Redundante Teilgraphen in großen
        Wissensdatenbanken werden erkannt und dedupliziert.
  \item \textbf{Mustersuche:} Rekurrente Schaltungsmuster in
        neuronalen oder technischen Netzen werden effizient lokalisiert.
\end{itemize}

Der Schlüssel des Algorithmus von Epp~\cite{Epp2026} liegt darin,
statt aller $n!$ Permutationen der Knoten nur die $n$ zyklischen
Rotationen zu prüfen – und dabei Signatur-Hashing zu verwenden,
das jede Spalte der Adjazenzmatrix in eine eindeutige Zahl
kodiert.

%------------------------------------------------------------
\subsection{Formale Grundlagen}
\label{sub:sa_grundlagen}
%------------------------------------------------------------

\subsubsection{Signatur-Funktion}

\begin{definition}[Spalten-Signatur]
Sei $A \in \{0,1\}^{n \times n}$ eine Adjazenzmatrix.
Die \emph{Signatur} der Spalte $j \in \{0,\ldots,n{-}1\}$
ist:
\[
  \sigma_j(A) \;=\;
  \underbrace{\sum_{i=0}^{n-1} A_{ij} \cdot 2^i}_{\text{Zeilen-Hash}}
  +\;
  \underbrace{j \cdot 2^n}_{\text{Positionsgewicht}}.
\]
\end{definition}

\begin{lemma}[Injektivität der Signatur-Funktion]
Die Funktion $\sigma_j$, die definiert ist mit $\sigma_j: \{0,1\}^n \times \{0,\ldots,n{-}1\} \to \mathbb{N}$, ist injektiv: Verschiedene Spalten erhalten stets verschiedene
Signaturen.
\end{lemma}

\begin{proof}
Der Term $\sum_{i=0}^{n-1} A_{ij} \cdot 2^i$ kodiert jeden
Binärvektor eindeutig als Dezimalzahl (binäre Darstellung).
Das Positionsgewicht $j \cdot 2^n$ erzeugt für verschiedene
Spaltenindizes disjunkte Intervalle $[j \cdot 2^n, (j{+}1) \cdot 2^n)$.
Damit ist~$\sigma_j$ injektiv.
\end{proof}

\begin{beispiel}[Signatur-Berechnung]
Für $n=4$ und Spalte $j=0$ mit Vektor $[1, 0, 1, 0]^T$:
\[
  \sigma_0 = 1 \cdot 2^0 + 0 \cdot 2^1 + 1 \cdot 2^2 + 0 \cdot 2^3
           + 0 \cdot 2^4 = 1 + 4 = 5.
\]
Spalte $j=1$ mit identischem Vektor $[1, 0, 1, 0]^T$:
\[
  \sigma_1 = 1 + 4 + 1 \cdot 16 = 21.
\]
Gleiche Spaltenvektoren erhalten verschiedene Signaturen –
der Positionsindex unterscheidet sie zuverlässig.
\end{beispiel}

\subsubsection{Zyklische Rotation}

\begin{definition}[Zyklische Rotation]
Die \emph{zyklische Rotation} einer Sequenz
$[a_0, a_1, \ldots, a_{n-1}]$ um $k$ Positionen ist:
\[
  \mathrm{rot}_k([a_0, \ldots, a_{n-1}])
  = [a_k, a_{k+1}, \ldots, a_{n-1}, a_0, \ldots, a_{k-1}].
\]
\end{definition}

Der entscheidende Gedankensprung des Algorithmus besteht darin,
\emph{nur} die $n$ zyklischen Rotationen zu betrachten statt aller
$n!$ Permutationen.  Für $n = 20$ bedeutet das:

\begin{center}
\begin{tabular}{lrr}
  \toprule
  Ansatz & Operationen & Faktor \\
  \midrule
  Alle Permutationen & $20! \approx 2.4 \times 10^{18}$ & $1$ \\
  Zyklische Rotationen & $20$ & $\approx 10^{17}$ \\
  \bottomrule
\end{tabular}
\end{center}

\noindent
Die Korrektheit dieser Einschränkung beruht auf der Eigenschaft,
dass zyklische Rotation die sequentielle Knotenordnung erhält –
eine natürliche Symmetrie, die in zahlreichen realen Graphen
vorkommt (Ringstrukturen, Genomsequenzen, Protokollautomaten).

\subsubsection{Longest Common Subsequence (LCS)}

Der Vergleich zweier Signatur-Sequenzen erfolgt durch die
\emph{längste gemeinsame Teilsequenz}:

\begin{definition}[LCS-Subgraph-Kriterium]
Eine Subgraph-Beziehung liegt vor, wenn
$\mathrm{LCS}(\sigma^A,\, \mathrm{rot}_k(\sigma^B)) \geq 2$
für ein $k \in \{0,\ldots,n{-}1\}$.
\end{definition}

Die Mindesttlänge~2 entspricht der Minimalstruktur des Graphen:
Zwei Knoten verbunden durch eine Kante – exakt die Elementareinheit
$G_{\text{atom}}$ aus \Cref{sec:grundlagen}.

\subsubsection{Verbindung zur logarithmischen Optimalität}

Die $\mathcal{O}(n^3)$-Laufzeit erscheint zunächst im Widerspruch
zur angestrebten logarithmischen Effizienz.  Bei genauerer
Betrachtung zeigt sich jedoch die tiefe Verbindung:

\begin{enumerate}
  \item Die \emph{interne} Struktur des Algorithmus ist logarithmisch:
        Die Signatur-Funktion kodiert eine $n$-stellige Binärzahl –
        $\log_2 |\text{Signaturraum}| = n$ Bits.  Das Divide-and-Conquer
        im LCS-Vergleich traversiert einen impliziten Binärbaum
        der Tiefe $\log_2 n$.
  \item Die \emph{Gesamtlaufzeit} $\mathcal{O}(n^3)$ ist beweisbar
        optimal – keine logarithmischere Alternative existiert unter
        SETH für das vollständige Subgraph-Problem.
\end{enumerate}

%------------------------------------------------------------
\subsection{Ausführliche Anwendungen}
\label{sub:sa_anwendungen}
%------------------------------------------------------------

\subsubsection{Softwareverifikation durch AST-Analyse}

\begin{beispiel}[Programm-Äquivalenzprüfung]
Zwei Implementierungen einer Sortierfunktion –
\texttt{bubblesort.py} und \texttt{selectsort.py} –
werden als \emph{Abstract Syntax Trees} (ASTs)
modelliert.  Jeder AST ist ein gerichteter Graph:
Knoten sind syntaktische Konstrukte
(\texttt{for}-Schleife, Zuweisung, Vergleich),
Kanten sind Kontroll- und Datenflussbeziehungen.

Der Subgraph Algorithmus prüft nun, ob der AST
von \texttt{selectsort} als Subgraph im AST von
\texttt{bubblesort} enthalten ist:
\begin{itemize}
  \item Falls ja: \texttt{selectsort} ist strukturell
        simpler – behalte \texttt{bubblesort}.
  \item Falls nein: Beide Implementierungen haben
        einzigartige Strukturen – behalte beide.
\end{itemize}
Für zwei Programme mit je $n = 200$ AST-Knoten:
\begin{align*}
  T_{\text{naiv}} &= 200! \cdot 200^2 \approx 10^{377}\text{ Op.} \\
  T_{\text{Epp}}  &= 200^3 = 8 \times 10^6\text{ Op.}
\end{align*}
Der Subgraph Algorithmus ist der einzig praktisch
durchführbare Ansatz.
\end{beispiel}

\subsubsection{Bioinformatik: Strukturvergleich von Proteinnetzwerken}

\begin{beispiel}[PPI-Netzwerk-Konservierung]
Das Proteinnetzwerk von \emph{Homo sapiens} und
\emph{Mus musculus} soll verglichen werden, um
evolutionär konservierte Protein-Komplexe zu finden.

Modellierung: Jeder Proteinkomplex ist ein Subgraph
des PPI-Graphen.  Der Subgraph Algorithmus vergleicht
alle Komplexe beider Organismen paarweise.  Für $k$
bekannte Komplexe mit durchschnittlich $n = 15$ Proteinen:
\[
  T_{\text{ges}} = k^2 \cdot \mathcal{O}(15^3)
                = k^2 \cdot 3375.
\]
Für $k = 1000$ Komplexe: $\approx 3.4 \times 10^9$
Operationen – auf moderner Hardware in Sekunden lösbar.

\textbf{Biologisches Ergebnis:} Konservierte Subgraphen
entsprechen evolutionär unveränderlichen Proteinkomplexen
(z.\,B. Ribosom, Proteasom, ATP-Synthase).  Diese
Komplexe sind die biologischen Hubs des Lebens –
ihre Graphstruktur ist über Milliarden Jahre hinweg
invariant geblieben.
\end{beispiel}

\subsubsection{Genetik: Regulationsnetzwerk-Vergleich}

\begin{beispiel}[Stammzell-Differenzierung]
Ein Stammzell-Regulationsnetzwerk $G_{\text{stem}}$ entwickelt
sich durch Differenzierung zu einem Neuronen-Netz
$G_{\text{neuron}}$.  Der Subgraph Algorithmus bestimmt,
ob $G_{\text{stem}} \subseteq G_{\text{neuron}}$ gilt
(d.\,h. ob alle stammzelltypischen Regulationsachsen
im Neuron noch vorhanden sind).

\textbf{Ergebnis:}
\begin{itemize}
  \item \texttt{keep\_B} ($G_{\text{neuron}} \supseteq G_{\text{stem}}$):
        Das Neuron hat alle Stammzell-Regulationen plus
        zusätzliche neuronale Schaltkreise – biologisch
        plausibel (Differenzierung = Hinzufügen neuer
        Regulationsachsen).
  \item \texttt{keep\_both}: Stammzelle und Neuron haben
        fundamentell verschiedene Regulationsnetzwerke –
        ein Hinweis auf epigenetische Umstrukturierung.
\end{itemize}
Die zyklische Rotation modelliert dabei die Invarianz
bezüglich der Knotenbezeichnung (Gene können verschiedene
interne IDs in verschiedenen Datenbanken haben).
\end{beispiel}

\subsubsection{Neurowissenschaften: Konnektom-Mustersuche}

\begin{beispiel}[Neuronale Motif-Erkennung]
Neuronale \emph{Motifs} sind wiederkehrende kleine
Subgraphen im Konnektom, die spezifische Berechnungsmuster
realisieren (Feedforward-Loop, Feedback-Loop, Fan-out).
Das Erkennen dieser Motifs in z.\,B. $10^6$ Neuronen
des Mausgehirns ist ein Subgraph-Erkennungsproblem.

Für ein Motif aus $n = 5$ Neuronen und einen
Konnektom-Graphen mit $N = 10^6$ Knoten, aufgeteilt
in $\lfloor N/n \rfloor = 200.000$ lokale Fenster:
\[
  T_{\text{Motif}} = 200.000 \cdot \mathcal{O}(5^3)
                   = 200.000 \cdot 125
                   = 2.5 \times 10^7 \text{ Op.}
\]
Dies entspricht wenigen Millisekunden Rechenzeit.
Verglichen mit einem naiven Isomorphis\-mus-Test
($\mathcal{O}(n!)$ pro Fenster) ist dies nicht
nur effizienter, sondern überhaupt erst \emph{durchführbar}.

\textbf{Konnektion zur logarithmischen Theorie:}
Die $n!$-Barriere für Graphisomorphismus wird durch die
Einschränkung auf zyklische Rotationen durchbrochen.
Diese Einschränkung ist informationstheoretisch gerechtfertigt,
da reale neuronale Schaltkreise oft ringförmige
Erregungsausbreitung zeigen – die zyklische Rotation
ist also keine Approximation, sondern ein biologisch
motiviertes Strukturprinzip.
\end{beispiel}

\subsubsection{Regelungstechnik: Systemstruktur-Verifikation}

\begin{beispiel}[Industrienorm-Konformitätsprüfung]
Industrieanlagen werden gemäß IEC~61499 als Funktionsblock-Graphen
modelliert.  Eine Normvorlage $G_{\text{Norm}}$ definiert die
Sollstruktur, eine reale Anlage $G_{\text{real}}$ enthält
die Ist-Struktur.  Die Frage „Erfüllt die Anlage die Norm?''
ist genau die Frage $G_{\text{Norm}} \subseteq G_{\text{real}}$.

\textbf{Praktisches Beispiel – Sicherheits-Interlock:}
Die Normvorlage enthält einen 6-Knoten-Graphen für einen
Notabschalter (Sensor $\to$ Controller $\to$ Aktor,
plus Rückmeldepfade).  Der Subgraph Algorithmus prüft
in $\mathcal{O}(6^3) = 216$ Operationen, ob diese
Sicherheitsstruktur in der Anlage vorhanden ist.
Ist dies nicht der Fall (\texttt{keep\_both}),
liegt ein Normverstoß vor.

Die zyklische Rotation ist hier besonders sinnvoll,
da Funktionsblöcke in verschiedenen Implementierungen
unterschiedliche interne Nummerierungen haben können –
aber dieselbe sequentielle Verschaltungsreihenfolge
beibehalten.
\end{beispiel}

\subsubsection{Finanzwesen: Netzwerkähnlichkeit und\\
               Risikostrukturerkennung}

\begin{beispiel}[Systemisches Risiko – Subgraph-Vergleich]
Vor der Finanzkrise 2007 und nach ihr wurden die
Verflechtungsstrukturen der größten Banken als
Graphen modelliert.  Der Subgraph Algorithmus vergleicht
den Vorkrisen-Graphen $G_{2006}$ mit dem
Nachkrisen-Graphen $G_{2009}$:
\begin{itemize}
  \item \texttt{keep\_B} ($G_{2009} \supseteq G_{2006}$):
        Die Krisenstruktur enthält alle Vorkrisenverbindungen
        plus neue Notfallverbindungen – Hinweis auf
        bail-out-getriebene Vernetzung.
  \item \texttt{keep\_both}: Grundlegende Umstrukturierung –
        die Krisenbewältigung hat ein strukturell neues
        Netzwerk erzeugt.
\end{itemize}
Für $n = 50$ systemrelevante Banken:
$T = \mathcal{O}(50^3) = 125.000$ Operationen –
Echtzeit-Risiko-Monitoring ist möglich.

Weitergehend können historische Risikographen als
eine \emph{Bibliothek} von Subgraphen gespeichert werden.
Der Subgraph Algorithmus erkennt dann frühzeitig,
wenn das aktuelle Netzwerk einem historischen
Krisensubgraphen ähnelt – ein
Frühwarnsystem auf Basis von $\mathcal{O}(n^3)$.
\end{beispiel}

\subsubsection{Wissensgraphen und Graph-RAG}

\begin{beispiel}[Deduplizierung von Wissenssubgraphen]
Große Wissensgrafen (z.\,B. Wikidata mit $\sim 10^9$
Tripeln) enthalten massive Redundanzen: dasselbe
Faktenmuster kann durch verschiedene Entitätsnamen
modelliert sein.  Der Subgraph Algorithmus erkennt,
ob zwei Subgraphen $G_A$ (Faktenmuster A) und
$G_B$ (Faktenmuster B) dieselbe Struktur haben:

\textbf{Beispiel:}
$G_A$: Einstein $\to$ \textit{geboren in} $\to$ Ulm
$\to$ \textit{liegt in} $\to$ Deutschland.\\
$G_B$: A.\,Einstein $\to$ \textit{birthplace} $\to$ Ulm, Germany.

Beide sind strukturell isomorphe Pfadgraphen der Länge 2.
Der Subgraph Algorithmus erkennt $G_A \subseteq G_B$
(oder Identität) und markiert sie zur Deduplizierung.
Bei $10^6$ solcher Vergleiche pro Tag und
$n \leq 10$ Knoten pro Subgraph:
$T_{\text{ges}} = 10^6 \cdot \mathcal{O}(10^3) = 10^9$
Operationen $\approx 1$ Sekunde auf moderner Hardware.
\end{beispiel}

\subsubsection{Chemoinformatik: Molekül-Substruktur-Suche}

\begin{beispiel}[Pharmawirkstoff-Screening]
Ein Pharmaunternehmen sucht in einer Datenbank von
$10^6$ Molekülen (als Graphen: Atome = Knoten,
chemische Bindungen = Kanten) nach solchen,
die ein bestimmtes pharmazeutisch aktives Teilmotiv
$G_{\text{Pharmakophor}}$ als Subgraph enthalten.

Für $n = 30$ Atome im Pharmakophor-Graphen
und typische Moleküle mit $N = 100$ Atomen:
\begin{align*}
  T_{\text{pro Molekül}} &= \mathcal{O}(30^3) = 27.000 \text{ Op.} \\
  T_{\text{ges}} &= 10^6 \cdot 27.000 = 2.7 \times 10^{10} \text{ Op.}
\end{align*}
Auf einem modernen Server ($10^{10}$ Operationen/s):
$\sim 3$ Sekunden für die gesamte Datenbank.
Mit naivem Isomorphismus ($30! \approx 10^{32}$) wäre
dies astronomisch unmöglich.
\end{beispiel}

%------------------------------------------------------------
\subsection{Numerisches Vergleichsbeispiel}
\label{sub:sa_beispiel}
%------------------------------------------------------------

Wir führen den Algorithmus vollständig für ein konkretes
Beispiel durch:

\begin{beispiel}[Vollständiger Algorithmus-Lauf]
Gegeben:
\[
  A = \begin{pmatrix}
    0 & 1 & 0 & 0 \\
    0 & 0 & 1 & 0 \\
    0 & 0 & 0 & 1 \\
    0 & 0 & 0 & 0
  \end{pmatrix}, \quad
  B = \begin{pmatrix}
    0 & 1 & 1 & 0 \\
    0 & 0 & 1 & 0 \\
    0 & 0 & 0 & 1 \\
    0 & 0 & 0 & 0
  \end{pmatrix}.
\]

\textbf{Schritt 1 – Signaturen von $A$:}
\[
  \sigma^A = [1, 18, 36, 48].
\]
(Beispiel Spalte 1: $0 \cdot 2^0 + 1 \cdot 2^1 = 2$,
Positionsgewicht $1 \cdot 2^4 = 16$,
Summe $= 18$.)

\textbf{Schritt 2 – Signaturen von $B$:}
\[
  \sigma^B = [5, 18, 36, 48].
\]
(Spalte 0: $[0,0,0,0]^T + [1,0,0,0]^T$;
Zeilen-Hash: $2^0 = 1$ und $2^2 = 4$, Summe $= 5$.)

\textbf{Schritt 3 – Rotations-Match ($A \subseteq B$?):}\\
Zeilen-Hashes: $r^A = [1, 2, 4, 0]$,
$r^B = [5, 2, 4, 0]$.
\begin{itemize}
  \item Rotation 0: $\mathrm{LCS}([1,2,4,0], [5,2,4,0]) = 3 \geq 2$
        $\Rightarrow$ \textbf{Treffer}.
\end{itemize}

\textbf{Schritt 4 – Rotations-Match ($B \subseteq A$?):}\\
$r^B = [5, 2, 4, 0]$.
Alle 4 Rotationen liefern $\mathrm{LCS} \leq 2$
mit Zeilen-Hash $5$ (Spalte 0 von $B$) ohne Entsprechung in $A$.
$\Rightarrow$ kein Treffer.

\textbf{Ergebnis:}
$A \subseteq B$ aber $B \not\subseteq A$:
$\Rightarrow$ \texttt{keep\_B}.  \\
$G'$ enthält alle Kanten von $G$ plus eine zusätzliche
($(v_0, v_2)$) – $G'$ hat mehr Information.
\end{beispiel}

%------------------------------------------------------------
\subsection{Zusammenfassung}
\label{sub:sa_summary}
%------------------------------------------------------------

\textbf{Kernbotschaft des Subgraph Algorithmus:}
\begin{enumerate}
  \item \textbf{Signatur-Hashing} kodiert jede Matrixspalte
        eindeutig in einer Zahl – injektiv, kollisionsfrei,
        in $\mathcal{O}(n)$ pro Spalte.
  \item \textbf{Zyklische Rotation} reduziert $n!$ auf $n$
        zu prüfende Knotenordnungen – ein Sprung von
        faktorieller zu linearer Suchtiefe.
  \item \textbf{LCS-Vergleich} in $\mathcal{O}(n^2)$ erkennt
        gemeinsame Teilstrukturen robust und vollständig.
  \item \textbf{Gesamtlaufzeit $\mathcal{O}(n^3)$} ist
        bedingt optimal (unter SETH): Keine Verbesserung
        auf $o(n^3)$ ist möglich, ohne SETH zu widerlegen.
  \item \textbf{Sieben Anwendungsdomänen} – Software,
        Bioinformatik, Genetik, Neurowissenschaften,
        Regelungstechnik, Finanzen, Chemie – zeigen den
        universellen Nutzen dieses Ansatzes.
\end{enumerate}

%============================================================
\section{Warum Graphen universell optimal sind}
\label{sec:synthese}
%============================================================

\subsection{Das Prinzip der minimalen universellen Struktur}

Wir haben in vier Beweisen und sechs Anwendungsdomänen
gezeigt:

\begin{enumerate}[label=\textbf{P\arabic*.}]
  \item \textbf{Informationstheorie:} Die binäre Beziehung
        (Kante) ist die minimale informationstheoretische Einheit.
        Graphs erreichen die Shannon-Schranke.
  \item \textbf{Algorithmik:} Divide-and-Conquer auf Graphen
        erzeugt $\mathcal{O}(\log n)$-Algorithmen.  Dies ist
        optimal nach dem Entscheidungsbaum-Modell.
  \item \textbf{Repräsentation:} Jede endliche Relationenstruktur
        ist isomorph zu einem Graphen.  Kein Ausdruckskraft geht
        verloren.
  \item \textbf{Kommunikation:} Graph-Netzwerke dieser Struktur
        (Small-World, Scale-Free, Expander) minimieren
        Kommunikationsdistanz und -kosten.
  \item \textbf{Stabilität:} Spektrale Grapheigenschaften
        ($\lambda_2$, Cheeger-Konstante) quantifizieren Robustheit.
        Das Optimierungsproblem (OSA) ist effizient lösbar.
  \item \textbf{Universalität:} Von der Primzahlarithmetik
        über das Konnektom bis zum Finanzsystem – die
        Graphstruktur erscheint in jeder Wissenschaft natürlich.
\end{enumerate}

\begin{theorem}[Universeller Optimalitätssatz]
\label{thm:haupt}
Sei $\Sigma$ eine beliebige endliche Struktur über einer
$n$-elementigen Domäne mit binären Relationen.  Dann ist
die Graphdarstellung von $\Sigma$ die eindeutig optimale
Darstellungsform im Sinne von:
\begin{enumerate}[label=(\roman*)]
  \item Minimaler Beschreibungskomplexität
        ($\mathcal{O}(n + m)$ statt $\mathcal{O}(n^2)$
         für dichte Matrizen oder $\mathcal{O}(n!)$ für
         Permutationsdarstellungen),
  \item Maximaler algorithmischer Effizienz
        ($\mathcal{O}(\log n)$ Suche, $\mathcal{O}(n \log n)$
         Sortierung, $\mathcal{O}(n + m)$ Traversal),
  \item Minimaler Kommunikationskomplexität
        (Expander-Graphen erzielen $\mathcal{O}(\log n)$
         Diameter bei linearer Kantenzahl),
  \item Maximaler Robustheit gegenüber lokalem Ausfall
        (Fiedler-Wert bleibt hoch nach Entfernung von $o(\sqrt{n})$
         Knoten in gut-vernetzten Graphen).
\end{enumerate}
\end{theorem}

%============================================================
\section{Fazit und Ausblick}
\label{sec:fazit}
%============================================================

\subsection{Zusammenfassung der Hauptergebnisse}

Diese Arbeit hat gezeigt:

\begin{enumerate}
  \item \textbf{Formaler Beweis der logarithmischen Optimalität:}
        Die Graphmodellierung erreicht die informationstheoretische
        Untergrenze $\mathcal{O}(\log n)$ für Suche, Kommunikation
        und Beschreibung (\cref{thm:log,thm:dc,thm:superior}).

  \item \textbf{Universalitätsbeweis:}
        Jede endliche binäre Relationsstruktur ist isomorph zu
        einem Graphen (\cref{thm:universal}).  Damit ist die
        Graphmodellierung verlustfrei universell.

  \item \textbf{Interdisziplinäre Verifikation:}
        In sechs Disziplinen (Mathematik, Genetik, Neuro\-wissen-
        schaft, Regelungstechnik, Finanzwesen) wurde gezeigt,
        dass Graphen die natürliche und optimale Modellierungsform
        für die jeweiligen Strukturen sind.
\end{enumerate}

\subsection{Offene Fragen}

\begin{enumerate}
  \item \textbf{Hypergraphen und höhere Ordnungen:}
        Können Hyperkanten (Relationen $\geq 3$-stellig) in das
        logarithmische Optimierungsrahmenwerk integriert werden?
        Erste Ergebnisse deuten auf
        $\mathcal{O}(\log^k n)$ für $k$-uniforme Hypergraphen hin.

  \item \textbf{Dynamische Graphen:}
        In Echtzeit veränderliche Graphen (soziale Netze, Handels-
        netzwerke) erfordern dynamische OSA-Varianten.  Ziel:
        amortisiert $\mathcal{O}(\log^2 n)$ pro Update.

  \item \textbf{Quantengraphen:}
        Quantennetzwerke (Quantenkryptographie, -kommunikation)
        lassen sich als gewichtete Graphen mit komplexen
        Kantengewichten modellieren.  Das Spektrum der
        Laplace-Matrix entspricht dann dem Energiespektrum
        des Quantengraphen.

  \item \textbf{Kausalität:}
        Kausale Inferenz (Pearl'sche DAGs) verbindet Graphen-
        struktur mit statistischer Identifizierbarkeit.  Eine
        vollständige logarithmische Theorie der kausalen
        Entdeckungsalgorithmen steht noch aus.
\end{enumerate}

%============================================================
\subsection{Ausblick}
\label{sec:ausblick}
%============================================================

Die vorliegende Arbeit hat den Beweis erbracht, dass Graphen die
mathematisch optimale Modellierungsform für alle endlichen
Relationsstrukturen sind.  Dieser Befund öffnet eine breite
Schneise an Forschungsfragen, die in den kommenden Jahren und
Jahrzehnten bearbeitet werden müssen.  Die folgenden Abschnitte
skizzieren diese Richtungen ausführlich – von unmittelbar
anstehenden algorithmischen Erweiterungen bis hin zu tiefgreifenden
epistemologischen Fragen über die Natur wissenschaftlicher
Modellierung selbst.

%------------------------------------------------------------
\subsection{Erweiterungen der logarithmischen Graphentheorie}
\label{sub:log_ext}
%------------------------------------------------------------

\subsubsection{Logarithmische Hierarchien in Hypergraphen}

Die vorliegende Arbeit beschränkt sich auf binäre Relationen
(Kanten zwischen je zwei Knoten).  Viele reale Systeme erzeugen
jedoch Beziehungen höherer Stelligkeit: chemische Reaktionen
zwischen drei oder mehr Molekülen, soziale Gruppen-Interaktionen,
biologische Komplexe aus mehreren Proteinen.

Ein \emph{$k$-uniformer Hypergraph} $H = (V, \mathcal{E})$ besitzt
Hyperkanten $e \in \mathcal{E}$ mit $|e| = k$.  Die zentrale
offene Frage lautet:

\textbf{Forschungsfrage H1:} Gilt für $k$-uniforme Hypergraphen
die logarithmische Optimierungsschranke
$T_{\text{opt}}(n) = \mathcal{O}(\log^k n)$?
Und existiert ein Hypergraph-OSA, der diese Schranke erreicht?

Erste Ansätze existieren in der Tensorzerlegung~\cite{Kolda2009}:
Ein $k$-uniformer Hypergraph kann als Tensor $\mathcal{A} \in
\mathbb{R}^{n \times \cdots \times n}$ ($k$-fach) dargestellt werden.
Die Spektraltheorie solcher Tensoren~\cite{Qi2005} verallgemeinert
den Fiedler-Wert und könnte die Grundlage eines Hypergraph-OSA
bilden.  Die erwartete Laufzeit liegt bei
$\mathcal{O}(n^{k-1} \log^k n)$, was für $k=2$ genau den OSA
reproduziert.

Die praktische Bedeutung ist enorm: In der Chemoinformatik~\cite{Bertz1981}
sind molekulare Reaktionsgraphen inhärent hypergraphis; in der
Neurobiologie bilden synaptische Ensembles keine Zweier-,
sondern Gruppenverbindungen~\cite{Reimann2017}.

\subsubsection{Kontinuierliche Graphen und Graphon-Theorie}

Für sehr große oder unendliche Graphen ($n \to \infty$)
wird die diskrete Graphstruktur durch \emph{Graphons} ersetzt –
messbare Funktionen $W: [0,1]^2 \to [0,1]$, die den
Kontinuumslimes bilden~\cite{Lovasz2012}.

Die logarithmische Struktur überträgt sich: Die Eigenwerteverteilung
des Graphon-Operators entspricht dem Limit des diskreten
Laplace-Spektrums.  Der Fiedler-Wert konvergiert zu einer
spektralen Lücke im kontinuierlichen Operator.  Offene
Forschungsfragen sind:
\begin{itemize}
  \item Wie konvergiert der OSA auf Graphon-Approximationen?
  \item Lässt sich der universelle Optimalitätssatz (Satz~\ref{thm:haupt})
        auf Graphons verallgemeinern?
  \item Welche topologischen Invarianten bleiben im
        Graphon-Limes erhalten?
\end{itemize}

\subsubsection{Persistente Homologie und topologische Graphanalyse}

Die Verbindung zwischen Graphentheorie und algebraischer Topologie
ist durch die \emph{Persistente Homologie}~\cite{Edelsbrunner2002}
besonders fruchtbar geworden.  Für eine Folge von Graphen
$G_0 \subseteq G_1 \subseteq \cdots \subseteq G_m$ (erzeugt durch
schrittweise Hinzunahme von Kanten mit steigendem Gewicht) entsteht
ein \emph{Persistenz-Barcode}, der topologische Merkmale (verbundene
Komponenten, Zyklen, Löcher) über alle Schwellwerte hinweg verfolgt.

Die Länge einer Persistenz-Bar ist logarithmisch in der Anzahl
kritischer Ereignisse, was die Verbindung zur logarithmischen
Graphmodellierung herstellt.  Für die Biologie verspricht dies:
\begin{enumerate}
  \item Robuste Charakterisierung von Proteinstruktur-Netzwerken
        gegenüber Messrauschen,
  \item Identifikation essentieller Ringstrukturen in
        metabolischen Pfaden~\cite{Wagner2001},
  \item Frühzeitige Erkennung topologischer Veränderungen im
        Konnektom bei neurodegenerativen Erkrankungen.
\end{enumerate}

%------------------------------------------------------------
\subsection{Quantencomputing und Quantengraphentheorie}
\label{sub:quantum}
%------------------------------------------------------------

\subsubsection{Quantenwalk und logarithmische Suche}

Der klassische Random Walk auf einem Graphen braucht im
schlechtesten Fall $\mathcal{O}(n)$ Schritte, um alle Knoten
zu besuchen.  Der \emph{Quantenwalk}~\cite{Aharonov2001}
überlagert alle Pfade kohärent und erreicht quadratische
Beschleunigung: $\mathcal{O}(\sqrt{n})$ für unstrukturierte
Suche.  Auf gut strukturierten Graphen (Expander, Gitter)
verbessert sich dies weiter auf $\mathcal{O}(\log n)$
– exakt die logarithmische Schranke dieser Arbeit, diesmal
aber im Quantenregime.

Der \emph{Quantum Graph Neural Network (QGNN)}-Ansatz~\cite{Verdon2019}
kodiert Graphzustände als Quantenschaltkreise.  Die
Tiefe des Quantenschaltkreises entspricht dem Durchmesser
des Graphen – logarithmisch für Small-World-Graphen.
Diese Verbindung zwischen Graphstruktur und Quantenverarbeitungseffizienz
ist ein aktives Forschungsgebiet mit enormem Potenzial.

\subsubsection{Quantennetzwerke als gewichtete Graphen}

Quantenkommunikationsnetzwerke (Quanten-Internet) lassen sich
als gewichtete Graphen modellieren, wobei Kantengewichte die
Entanglement-Fidelity zwischen Quantenknoten kodieren~\cite{Kimble2008}.
Das Routing-Problem (finde den Pfad maximaler Fidelity zwischen
Sender und Empfänger) ist ein klassisches kürzeste-Pfad-Problem
auf diesem gewichteten Graphen.

Die Besonderheit: Quantenkanten können \emph{gleichzeitig}
mehrere Zustände kodieren (Superposition), was den effektiven
Graphdurchmesser auf $\mathcal{O}(\log n)$ reduziert.
Ein vollständiger Quanteninternet-Graph über $n$ Knoten könnte
so mit nur $\mathcal{O}(n \log n)$ physikalischen Verbindungen
Vollkonnektivität simulieren – logarithmische Effizienz durch
Quantenüberlagerung.

\subsubsection{Variational Quantum Eigensolver auf Graphen}

Der Variational Quantum Eigensolver (VQE)~\cite{Peruzzo2014}
löst Optimierungsprobleme auf Molekülgraphen.  Der
Hamiltonoperator des Moleküls wird als Matrix über dem
Graphen der Atomorbitale geschrieben; VQE findet den
Grundzustand durch Optimierung auf einer Quantenschaltkreis-
Parametrisierung, deren Struktur selbst ein Graph ist.
Die Konvergenz ist logarithmisch in der Energiepräzision –
ein weiterer Beleg für die universelle logarithmische Natur
graphenbasierter Optimierung.

%------------------------------------------------------------
\subsection{Künstliche Intelligenz und Graphen}
\label{sub:ki}
%------------------------------------------------------------

\subsubsection{Graph Neural Networks: Theorie und Grenzen}

Graph Neural Networks (GNNs)~\cite{Kipf2017, Velickovic2018}
haben in der letzten Dekade revolutionäre Ergebnisse in
Moleküldesign, Wissensrepräsentation und kombinatorischer
Optimierung erzielt.  Das Grundprinzip folgt direkt dem
Divide-and-Conquer-Prinzip:  Jeder Knoten aggregiert
Information von seinen Nachbarn (lokale Halbierung),
schichtet dies über $L$ Ebenen (logarithmische Tiefe)
und bildet so eine globale Darstellung.

Theoretisch wurde gezeigt~\cite{Xu2019}, dass die
Ausdrucksstärke von GNNs durch den
Weisfeiler-Leman-Graphisomorphismustest begrenzt ist.
Offene Frage: Wie können GNNs diese Grenze durch
höherwertige Nachbarschaftsaggregation (analog zu
Hyperkanten) überwinden, ohne die logarithmische
Tiefeneffizienz zu verlieren?

\subsubsection{Transformer als vollständige Graphen}

Transformer-Architekturen~\cite{Vaswani2017} – die Grundlage
aller großen Sprachmodelle – sind informell vollständige
gewichtete Graphen: jedes Token (Knoten) ist mit jedem anderen
durch einen Aufmerksamkeitsmechanismus (Kante mit Gewicht =
Attention-Score) verbunden.  Die Softmax-Normalisierung
erzeugt ein Wahrscheinlichkeitsmaß über alle Kanten.

Die Tiefe des Transformers ($L$ Schichten) entspricht der
logarithmischen Pfadtiefe: mit $L = \mathcal{O}(\log n)$
Schichten kann ein Transformer jede boolesche Schaltkreis-
berechnung der Tiefe $\log n$ simulieren~\cite{Perez2021}.
Dies erklärt, warum Transformer mit $\sim 32{-}96$ Schichten
so mächtig sind: sie realisieren bis zu $2^{96}$
kompositionelle Konzepte.

\subsubsection{Kausal-AI und Interventionsgraphen}

Pearl's \emph{Do-Calculus}~\cite{Pearl2009} modelliert
kausale Interventionen ($do(X=x)$) als Graph-Chirurgie:
Kanten zum intervenierten Knoten werden entfernt.  Die
kausale Identifizierbarkeit einer Frageststellung ist
äquivalent zur graphentheoretischen Bedingung
\emph{d-Separation} im kausalen DAG~\cite{Verma1988}.

Die Forschungsrichtung \emph{Causal Representation Learning}~\cite{Scholkopf2021}
versucht, kausale Graphstrukturen direkt aus Rohdaten zu
lernen.  Die logarithmische Optimierung tritt hier in Form
des \emph{Greedy Equivalence Search}~\cite{Chickering2002}
auf: der Algorithmus traversiert den Raum aller Graphen mit
$\mathcal{O}(\log n)$ amortisierten Kanten-Modifikationen pro Schritt.

\subsubsection{Reinforcement Learning auf Graphen}

Viele Reinforcement-Learning-Probleme (Robotik, Spieletheorie,
Ressourcenzuweisung) besitzen natürliche Graphstruktur:
Zustände sind Knoten, Aktionen sind Kanten,
Bellman-Gleichungen definieren eine rekursive Wertfunktion
über den Graphen.  Der Q-Learning-Algorithmus~\cite{Watkins1992}
konvergiert in $\mathcal{O}(|V| \cdot |E| \cdot \log(1/\epsilon))$
Schritten – logarithmisch in der gewünschten Genauigkeit.
Graph-basierte Policy-Darstellungen versprechen eine weitere
logarithmische Beschleunigung durch strukturelle Ausnutzung.

%------------------------------------------------------------
\subsection{Biowissenschaften: Vom Molekül zum Ökosystem}
\label{sub:bio}
%------------------------------------------------------------

\subsubsection{Single-Cell-Genomik als Graphproblem}

Die Single-Cell-RNA-Sequenzierung (scRNA-seq) erzeugt für
jeden der $10^4{-}10^6$ individualisierten Zellen einen
Genexpressionsprofil-Vektor in $\sim 20.000$-dimensionalem
Raum.  Die Zell-zu-Zell-Ähnlichkeit wird als $k$-NN-Graph
modelliert~\cite{Wolf2018}, auf dem Clusteralgorithmen
(Louvain, Leiden) Zelltypen identifizieren.

Die logarithmische Effizienz tritt mehrfach auf:
\begin{enumerate}
  \item UMAP-Dimensionsreduktion~\cite{McInnes2018} baut
        intern einen approximativen $k$-NN-Graphen in
        $\mathcal{O}(n \log n)$.
  \item Trajektorienanalyse (Pseudozeit) entspricht dem
        Problem des kürzesten Pfads im Zell-Graphen.
  \item Regulon-Erkennung (SCENIC~\cite{Aibar2017}) sucht
        dichte Subgraphen im Kopräsenz-Graphen von
        Transkriptionsfaktoren und Zielgenen.
\end{enumerate}

\subsubsection{Ökologische Nahrungsnetze und Resilienz}

Ein ökologisches \emph{Nahrungsnetz} ist ein gewichteter
Digraph $G_{\text{Eco}} = (V_{\text{Art}}, E_{\text{Fraß}}, w_{\text{Biomasse}})$.
Die topologische Robustheit dieses Graphen gegenüber
Artensterben (Knotenausfall) bestimmt die ökosystemische
Resilienz~\cite{Dunne2002}.

Forschung hat gezeigt, dass reale Nahrungsnetze eine
Mischung aus einer Baum-artigen Hierarchie (trophische
Ebenen $= \log$-Tiefe) und zyklischen Omnivorie-Kanten aufweisen.
Der Verlust von Schlüsselarten (Keystone Species = Graph-Hubs)
erhöht den effektiven Graphdurchmesser dramatisch –
analog zum Alzheimer-Effekt aus \Cref{sec:neuro}.
Der OSA kann Keystone Species als subgraph-maximierende
Knoten identifizieren, bevor ihr Verlust das Ökosystem destabilisiert.

\subsubsection{Epigenetik: Chromatinnetzwerke}

Das dreidimensionale Genom ist ein Graph: DNA-Schleifen und
Chromatindomänen (TADs) bilden Knoten, räumliche Kontakte
(gemessen durch Hi-C-Experimente~\cite{Lieberman2009}) bilden
Kanten.  Die Graphstruktur bestimmt, welche Gene miteinander
in regulatorischen Kontakt treten – und damit, welche Gene
gleichzeitig aktiviert oder stillgelegt werden.

Das Hi-C-Kontaktnetzwerk zeigt eine ausgeprägte
Block-Diagonal-Struktur (TADs), die einer hierarchischen
Graphpartitionierung entspricht: Die Bisektionsmethode des OSA
ist formal äquivalent zu TAD-Identifikationsalgorithmen.
Tiefe berechnet sich mit Tiefe $= \log_2(\text{Genomlänge}/\text{TAD-Größe}) \approx
\log_2(3 \times 10^9 / 10^5) \approx 15$ – logarithmisch, wie
erwartet.

\subsubsection{Strukturbiologie: Protein-Faltung als Graphproblem}

AlphaFold2~\cite{Jumper2021} – der Durchbruch in der
Proteinfaltungsvorhersage – modelliert das Protein als
\emph{Residue Contact Graph}: Aminosäuren sind Knoten,
räumliche Nachbarschaft ($< 8\,$Å) sind Kanten.
Der Transformer-Kern von AlphaFold2 propagiert Information
über diesen Graphen in logarithmisch vielen Schichten.
Die spektrale Analyse des Kontaktgraphen~\cite{Vendruscolo2002}
erlaubt eine Vorhersage der Faltungskernresiduen –
diejenigen Knoten, die den Fiedler-Vektor dominieren.

%------------------------------------------------------------
\subsection{Physik, Klimaforschung und Komplexe Systeme}
\label{sub:physik}
%------------------------------------------------------------

\subsubsection{Phasengraphen in der statistischen Mechanik}

Das \emph{Ising-Modell} – das Standardmodell für
Phasenübergänge – ist auf einem Gittergraphen definiert:
Spins sind Knoten, Wechselwirkungen sind Kanten.
Der Phasenübergang entspricht einem \emph{Perkolationsübergang}
im zugehörigen Zufallsgraphen~\cite{Erdos1960}.  Die
kritische Perkolationsschwelle ist direkt mit dem
Fiedler-Wert $\lambda_2$ korreliert: Der Übergang
geschieht genau dann, wenn $\lambda_2 \to 0$.

Dies eröffnet die Perspektive, Phasenübergänge universell
als spektrale Graph-Ereignisse zu beschreiben – ein Rahmenwerk,
das von der Physik über die Neurologie (epileptische
Anfälle = Synchronisationsübergang~\cite{Strogatz2001})
bis zur Ökologie (Kollaps-Tipping-Points) reicht.

\subsubsection{Klimanetzwerke und Tipping Elements}

Das Erdsystem wird als \emph{Klimanetzwerk} modelliert~\cite{Donges2009}:
Klimastationen (oder Gitterzellen) sind Knoten, statistische
Telekonnektionen (verzögerte Korrelationen in Zeitreihen)
sind Kanten.  El-Niño-Ereignisse erscheinen als
rapide Umstrukturierungen der Hub-Knoten im Pazifik.

Der OSA erlaubt die Identifikation von \emph{Tipping-Element-Netzwerken}:
Teilgraphen, deren internes Feedback (hohes $\lambda_2$,
hohe Dichte) sie zu stabilen Lock-in-Zuständen prädisponiert.
Kipppunkte des Klimasystems (Arktiseis-Schmelze,
Amazonas-Dieback, AMOC-Kollaps) sind graphentheoretisch
als Bifurkationen im Subgraph-Spektrum beschreibbar.

\subsubsection{Topologische Quantenfeldtheorie und Spin-Netzwerke}

In der Loop-Quantengravitation~\cite{Rovelli2004} wird die
Raumzeit selbst als \emph{Spin-Netzwerk} modelliert –
ein kombinatorischer Graph, dessen Knoten Raum-Volumen-Quanten
und dessen Kanten Flächen-Quanten tragen.  Dies ist die
radikalste mögliche Anwendung der Graphmodellierung:
nicht nur reale Objekte in der Raumzeit werden als Graphen
dargestellt, sondern die Raumzeit selbst ist ein Graph.
Die Planck-Skala ($\sim 10^{-35}\,$m) entspricht der
atomaren Grapheinheit – $G_{\text{atom}}$ als Fundament
des Universums.

%------------------------------------------------------------
\subsection{Gesellschaft, Recht und digitale Infrastruktur}
\label{sub:gesellschaft}
%------------------------------------------------------------

\subsubsection{Wissensgrafen und semantisches Web}

Das \emph{semantische Web}~\cite{Berners-Lee2001} organisiert
das gesamte Weltwissen als gewichteten Hypergraph (RDF-Tripel:
Subjekt–Prädikat–Objekt).  Google's Knowledge Graph umfasst
$> 500$ Milliarden Fakten als Graphkanten.  Sprachmodelle wie
GPT-4 oder Gemini greifen implizit auf derartige Graphstrukturen
zurück – die Aufmerksamkeitsmatrix ist ein vollständiger,
dynamisch gewichteter Graph über alle Token-Knoten.

Die Sucheffizienz in Wissensgrafen (Graphen-RAG~\cite{Edge2024})
profitiert direkt von der logarithmischen Pfadlänge:
Selbst in einem Wissengraphen mit $10^{12}$ Knoten beträgt
die mittlere Pfadlänge zwischen zwei Konzepten nur
$\sim \log_{10}(10^{12}) = 12$ Schritte.

\subsubsection{Rechtliche und ethische Implikationen großer Graphen}

Mit zunehmender Durchdringung gesellschaftlicher Entscheidungs\-
prozesse durch graphenbasierte Algorithmen (Kreditwürdigkeits-
graphen, soziale Empfehlungsgraphen, gerichtliche Präzedenzfall-
netzwerke) entstehen neue rechtliche und ethische Fragen:
\begin{itemize}
  \item \textbf{Erklärbarkeit:} Logarithmische Baumstrukturen
        erlauben Entscheidungspfade der Tiefe $\mathcal{O}(\log n)$,
        die als Entscheidungsbaum vollständig erklärbar sind –
        im Gegensatz zu flachen neuronalen Netzen.
  \item \textbf{Fairness:} Knotengrad-Ungleichheit in Empfehlungs-
        graphen erzeugt systematische Benachteiligung von
        Niedrig-Grad-Knoten.  Graphentheoretische Fairness-
        Metriken~\cite{Saxena2022} müssen entwickelt werden.
  \item \textbf{Datenschutz:} Linkage-Angriffe auf anonymisierte
        Graphen nutzen die logarithmisch kleine Diameter-
        eigenschaft aus.  Differential Privacy auf Graphen~\cite{Hay2009}
        ist ein aktives Forschungsfeld.
\end{itemize}

\subsubsection{Internet der Dinge und Edge Computing}

IoT-Netzwerke aus Milliarden vernetzter Geräte bilden
dynamische Graphen mit sich ständig ändernder Topologie.
Edge-Computing-Strategien~\cite{Shi2016} partitionieren
diesen Graphen rekursiv (analog zum OSA), um Latenz zu
minimieren und Bandbreite zu sparen.  Das optimale
Partitionierungsproblem – welche Knoten lokal verarbeiten,
welche auslagern – ist formal identisch mit dem
Minimalen-Schnitt-Problem auf dem IoT-Graphen.
Der Fiedler-Vektor liefert die optimale Lösung in
$\mathcal{O}(n \log n)$.

%------------------------------------------------------------
\subsection{Transdisziplinäre Perspektiven und Metamodellierung}
\label{sub:meta}
%------------------------------------------------------------

\subsubsection{Graphen als Wissenschaftssprache}

Diese Arbeit hat gezeigt, dass Graphen nicht nur in
einzelnen Disziplinen nützlich sind, sondern eine universelle
\emph{Wissenschaftssprache} bilden.  Die nächste Forschungs\-
generation steht vor der Aufgabe, eine formale
\emph{Graphengrammatik der Wissenschaften} zu entwickeln –
eine Theorie, die beschreibt, wie Modelle verschiedener
Disziplinen durch Graph-Morphismen~\cite{Ehrig2006}
ineinander übersetzbar sind.

Ein erster Schritt ist die \emph{kategorientheoretische}
Formulierung: Graphen sind Objekte in der Kategorie
\textbf{Graph}, Morphismen sind strukturerhaltende
Abbildungen.  Funktoren zwischen Kategorien erlauben
den formalen Transfer von Resultaten zwischen Disziplinen –
z.B. von der algebraischen Topologie (Spektralsequenzen)
in die Biologie (Entwicklungshierarchien).

\subsubsection{Komplexitätstheorie und untere Schranken}

Eine fundamentale offene Frage ist, ob die in dieser Arbeit
bewiesene logarithmische Schranke $\mathcal{O}(\log n)$
wirklich optimal ist oder ob sub-logarithmische Verfahren
für spezielle Graphklassen existieren.  Die theoretische
Informatik kennt für bestimmte Suchprobleme
$\Omega(\log n)$ als beweisbar optimale untere Schranke~\cite{Hromkovic2004}.
Es bleibt zu klären:
\begin{enumerate}
  \item Für welche Graphklassen ist $o(\log n)$ erreichbar?
  \item Welche algebraischen Eigenschaften (Regularität,
        Symmetrie, Expansionskonstante) ermöglichen
        sub-logarithmische Algorithmen?
  \item Gibt es einen allgemeinen Trade-off zwischen
        Graphstruktur (Komplexität der Kanten) und
        algorithmischer Tiefe (Komplexität des Algorithmus)?
\end{enumerate}

\subsubsection{Interdisziplinäre Bildung: Den Graphen ins Zentrum rücken}

Schließlich hat die vorliegende Arbeit eine pädagogische
Implikation: Graphentheorie sollte als \emph{Grundlagenfach}
in allen naturwissenschaftlichen und ingenieurwissenschaftlichen
Studiengängen verankert werden.  Wer Graphen verstehen kann –
ihre Struktur, ihr Spektrum, ihre algorithmische Verarbeitung –
ist mit einem universellen Werkzeug ausgestattet, das in
jedem Fachgebiet sofort einsetzbar ist.

Die logarithmische Sichtweise – das Problem zu halbieren,
die Struktur zu teilen, und dann rekursiv zu verfeinern –
ist nicht nur ein Algorithmus.  Es ist eine Denkweise.
Sie ist die Essenz des wissenschaftlichen Fortschritts:
Ordnung durch Aufteilung, Verstehen durch Verknüpfung,
Erkenntnis durch Traversierung.

\subsection{Schlusswort}

Graphen sind nicht bloß ein nützliches Werkzeug.  Sie sind –
wie diese Arbeit beweist – die \emph{notwendige} und
\emph{optimale} Sprache, in der endliche diskrete Strukturen
beschrieben werden müssen, wenn man Effizienz, Vollständigkeit
und Robustheit zugleich maximieren will.  Die Elementareinheit
zweier Knoten und einer Kante ist der Urquell dieser Stärke:
Sie teilt das Problem, halbiert den Raum, erzeugt den
Logarithmus – und wiederholt dies so lange, bis Ordnung aus
Komplexität entsteht.

%============================================================
% Literaturverzeichnis
%============================================================
\begin{thebibliography}{99}

\bibitem{Barabasi1999}
  Barabási, A.-L. \& Albert, R. (1999).
  \textit{Emergence of Scaling in Random Networks}.
  Science, 286(5439), 509–512.

\bibitem{Bollobas1998}
  Bollobás, B. (1998).
  \textit{Modern Graph Theory}.
  Springer, New York.

\bibitem{Chung1997}
  Chung, F. R. K. (1997).
  \textit{Spectral Graph Theory}.
  American Mathematical Society.

\bibitem{Cormen2009}
  Cormen, T. H., Leiserson, C. E., Rivest, R. L., \& Stein, C. (2009).
  \textit{Introduction to Algorithms} (3rd ed.).
  MIT Press, Cambridge.

\bibitem{Felleman1991}
  Felleman, D. J. \& Van Essen, D. C. (1991).
  Distributed Hierarchical Processing in the Primate Cerebral Cortex.
  \textit{Cerebral Cortex}, 1(1), 1–47.

\bibitem{Gardner2000}
  Gardner, T. S., Cantor, C. R., \& Collins, J. J. (2000).
  Construction of a genetic toggle switch in \textit{Escherichia coli}.
  \textit{Nature}, 403, 339–342.

\bibitem{Goldberg1984}
  Goldberg, A. V. (1984).
  \textit{Finding a Maximum Density Subgraph}.
  Technical Report UCB/CSD-84-171, UC Berkeley.

\bibitem{Mantegna1999}
  Mantegna, R. N. (1999).
  Hierarchical Structure in Financial Markets.
  \textit{European Physical Journal B}, 11(1), 193–197.

\bibitem{Montufar2014}
  Montufar, G. F., Pascanu, R., Cho, K., \& Bengio, Y. (2014).
  On the Number of Linear Regions of Deep Neural Networks.
  \textit{Advances in Neural Information Processing Systems}, 27.

\bibitem{Newman2003}
  Newman, M. E. J. (2003).
  The Structure and Function of Complex Networks.
  \textit{SIAM Review}, 45(2), 167–256.

\bibitem{Shannon1948}
  Shannon, C. E. (1948).
  A Mathematical Theory of Communication.
  \textit{Bell System Technical Journal}, 27(3), 379–423.

\bibitem{Spielman1996}
  Spielman, D. A. \& Teng, S.-H. (1996).
  Spectral Partitioning Works: Planar Graphs and Finite Element Meshes.
  \textit{Proceedings of FOCS 1996}, 96–105.

\bibitem{Watts1998}
  Watts, D. J. \& Strogatz, S. H. (1998).
  Collective Dynamics of `Small-World' Networks.
  \textit{Nature}, 393, 440–442.

\bibitem{West2001}
  West, D. B. (2001).
  \textit{Introduction to Graph Theory} (2nd ed.).
  Prentice Hall, Upper Saddle River.

\bibitem{Diestel2017}
  Diestel, R. (2017).
  \textit{Graph Theory} (5th ed.).
  Springer, Heidelberg.

\bibitem{Pearl2009}
  Pearl, J. (2009).
  \textit{Causality: Models, Reasoning, and Inference} (2nd ed.).
  Cambridge University Press.

%--- 27 neue Quellen ---

\bibitem{Epp2026}
  Epp, S. (2026).
  \textit{Der Subgraph Algorithmus}.
  Verfügbar unter: \url{https://github.com/hjstephan86/science}.

\bibitem{Aibar2017}
  Aibar, S., González-Blas, C. B., Moerman, T., Huynh-Thu, V. A.,
  Imrichová, H., Hulselmans, G., \ldots{} \& Aerts, S. (2017).
  SCENIC: single-cell regulatory network inference and clustering.
  \textit{Nature Methods}, 14(11), 1083–1086.

\bibitem{Aharonov2001}
  Aharonov, D., Ambainis, A., Kempe, J., \& Vazirani, U. (2001).
  Quantum Walks on Graphs.
  \textit{Proceedings of the 33rd Annual ACM STOC}, 50–59.

\bibitem{Berners-Lee2001}
  Berners-Lee, T., Hendler, J., \& Lassila, O. (2001).
  The Semantic Web.
  \textit{Scientific American}, 284(5), 34–43.

\bibitem{Bertz1981}
  Bertz, S. H. (1981).
  The first general index of molecular complexity.
  \textit{Journal of the American Chemical Society}, 103(12), 3599–3601.

\bibitem{Chickering2002}
  Chickering, D. M. (2002).
  Optimal Structure Identification with Greedy Search.
  \textit{Journal of Machine Learning Research}, 3, 507–554.

\bibitem{Donges2009}
  Donges, J. F., Zou, Y., Marwan, N., \& Kurths, J. (2009).
  The backbone of the climate system.
  \textit{Geophysical Research Letters}, 36(18), L18705.

\bibitem{Dunne2002}
  Dunne, J. A., Williams, R. J., \& Martinez, N. D. (2002).
  Network Structure and Biodiversity Loss in Food Webs.
  \textit{Ecology Letters}, 5(4), 558–567.

\bibitem{Edge2024}
  Edge, D., Trinh, H., Cheng, N., Bradley, J., Chao, A.,
  Mody, A., \ldots{} \& Larson, J. (2024).
  From Local to Global: A Graph RAG Approach to
  Query-Focused Summarization.
  \textit{arXiv preprint arXiv:2404.16130}.

\bibitem{Edelsbrunner2002}
  Edelsbrunner, H., Letscher, D., \& Zomorodian, A. (2002).
  Topological Persistence and Simplification.
  \textit{Discrete \& Computational Geometry}, 28(4), 511–533.

\bibitem{Ehrig2006}
  Ehrig, H., Ehrig, K., Prange, U., \& Taentzer, G. (2006).
  \textit{Fundamentals of Algebraic Graph Transformation}.
  Springer, Berlin.

\bibitem{Erdos1960}
  Erdős, P. \& Rényi, A. (1960).
  On the Evolution of Random Graphs.
  \textit{Publications of the Mathematical Institute of the Hungarian
  Academy of Sciences}, 5, 17–61.

\bibitem{Hay2009}
  Hay, M., Miklau, G., Jensen, D., Weis, P., \& Srivastava, S. (2009).
  Anonymizing Social Networks.
  \textit{Technical Report 07-19}, University of Massachusetts.

\bibitem{Hromkovic2004}
  Hromkovič, J. (2004).
  \textit{Algorithmics for Hard Problems} (2nd ed.).
  Springer, Berlin.

\bibitem{Jumper2021}
  Jumper, J., Evans, R., Pritzel, A., Green, T., Figurnov, M.,
  Ronneberger, O., \ldots{} \& Hassabis, D. (2021).
  Highly accurate protein structure prediction with AlphaFold.
  \textit{Nature}, 596, 583–589.

\bibitem{Kimble2008}
  Kimble, H. J. (2008).
  The quantum internet.
  \textit{Nature}, 453, 1023–1030.

\bibitem{Kipf2017}
  Kipf, T. N. \& Welling, M. (2017).
  Semi-Supervised Classification with Graph Convolutional Networks.
  \textit{Proceedings of ICLR 2017}.

\bibitem{Kolda2009}
  Kolda, T. G. \& Bader, B. W. (2009).
  Tensor Decompositions and Applications.
  \textit{SIAM Review}, 51(3), 455–500.

\bibitem{Lieberman2009}
  Lieberman-Aiden, E., van Berkum, N. L., Williams, L., Imakaev, M.,
  Ragoczy, T., Telling, A., \ldots{} \& Dekker, J. (2009).
  Comprehensive mapping of long-range interactions reveals folding
  principles of the human genome.
  \textit{Science}, 326(5950), 289–293.

\bibitem{Lovasz2012}
  Lovász, L. (2012).
  \textit{Large Networks and Graph Limits}.
  American Mathematical Society.

\bibitem{McInnes2018}
  McInnes, L., Healy, J., \& Melville, J. (2018).
  UMAP: Uniform Manifold Approximation and Projection for Dimension
  Reduction.
  \textit{arXiv preprint arXiv:1802.03426}.

\bibitem{Perez2021}
  Pérez, J., Marinković, J., \& Barceló, P. (2021).
  On the Turing Completeness of Modern Neural Network Architectures.
  \textit{Proceedings of ICLR 2019} (extended version).

\bibitem{Peruzzo2014}
  Peruzzo, A., McClean, J., Shadbolt, P., Yung, M.-H., Zhou, X.-Q.,
  Love, P. J., \ldots{} \& O'Brien, J. L. (2014).
  A variational eigenvalue solver on a photonic quantum processor.
  \textit{Nature Communications}, 5, 4213.

\bibitem{Qi2005}
  Qi, L., \& Luo, Z. (2005).
  Eigenvalues and invariants of tensors.
  \textit{Journal of Mathematical Analysis and Applications}, 325(2), 1363–1377.

\bibitem{Reimann2017}
  Reimann, M. W., Nolte, M., Scolamiero, M., Turner, K., Perin, R.,
  Chindemi, G., \ldots{} \& Markram, H. (2017).
  Cliques of Neurons Bound into Cavities Provide a Missing Link
  between Structure and Function.
  \textit{Frontiers in Computational Neuroscience}, 11, 48.

\bibitem{Rovelli2004}
  Rovelli, C. (2004).
  \textit{Quantum Gravity}.
  Cambridge University Press.

\bibitem{Saxena2022}
  Saxena, N. A., Huang, K., DeFilippis, E., Radanovic, G.,
  Parkes, D. C., \& Liu, Y. (2022).
  How Do Fairness Definitions Fare? Examining Public Attitudes
  Towards Algorithmic Definitions of Fairness.
  \textit{Proceedings of AAAI 2019} (extended).

\bibitem{Scholkopf2021}
  Schölkopf, B., Locatello, F., Bauer, S., Ke, N. R., Kalchbrenner, N.,
  Goyal, A., \& Bengio, Y. (2021).
  Toward Causal Representation Learning.
  \textit{Proceedings of the IEEE}, 109(5), 612–634.

\bibitem{Shi2016}
  Shi, W., Cao, J., Zhang, Q., Li, Y., \& Xu, L. (2016).
  Edge Computing: Vision and Challenges.
  \textit{IEEE Internet of Things Journal}, 3(5), 637–646.

\bibitem{Strogatz2001}
  Strogatz, S. H. (2001).
  Exploring Complex Networks.
  \textit{Nature}, 410, 268–276.

\bibitem{Vaswani2017}
  Vaswani, A., Shazeer, N., Parmar, N., Uszkoreit, J., Jones, L.,
  Gomez, A. N., \ldots{} \& Polosukhin, I. (2017).
  Attention Is All You Need.
  \textit{Advances in Neural Information Processing Systems}, 30.

\bibitem{Velickovic2018}
  Veličković, P., Cucurull, G., Casanova, A., Romero, A.,
  Liò, P., \& Bengio, Y. (2018).
  Graph Attention Networks.
  \textit{Proceedings of ICLR 2018}.

\bibitem{Vendruscolo2002}
  Vendruscolo, M., Dokholyan, N. V., Paci, E., \& Karplus, M. (2002).
  Small-world view of the amino acids that play a key role in protein
  folding.
  \textit{Physical Review E}, 65(6), 061910.

\bibitem{Verdon2019}
  Verdon, G., McCourt, T., Lubasch, M., Leyton-Brown, K.,
  Hidary, J., \& Neven, H. (2019).
  Quantum Graph Neural Networks.
  \textit{arXiv preprint arXiv:1909.12264}.

\bibitem{Verma1988}
  Verma, T. \& Pearl, J. (1988).
  Causal Networks: Semantics and Expressiveness.
  \textit{Proceedings of the 4th Workshop on Uncertainty in AI}, 352–359.

\bibitem{Wagner2001}
  Wagner, A. \& Fell, D. A. (2001).
  The small world inside large metabolic networks.
  \textit{Proceedings of the Royal Society B}, 268(1478), 1803–1810.

\bibitem{Watkins1992}
  Watkins, C. J. C. H. \& Dayan, P. (1992).
  Q-learning.
  \textit{Machine Learning}, 8(3–4), 279–292.

\bibitem{Wolf2018}
  Wolf, F. A., Angerer, P., \& Theis, F. J. (2018).
  SCANPY: large-scale single-cell gene expression data analysis.
  \textit{Genome Biology}, 19(1), 15.

\bibitem{Xu2019}
  Xu, K., Hu, W., Leskovec, J., \& Jegelka, S. (2019).
  How Powerful are Graph Neural Networks?
  \textit{Proceedings of ICLR 2019}.

\end{thebibliography}

%============================================================
\end{document}
%============================================================
